\chapter{Test configurations, norms, and filtrations}
\index{test configuration}\index{filtration}\index{non-Archimedean norm}
\label{chap:bj-test-configurations}

Test configurations are one-parameter degenerations.  Three related
encodings are useful for different purposes; their equivalence requires
the polarization and closure hypotheses stated below.  The total space is
best for
intersection theory, the Rees algebra is best for finite generation, and the
induced filtration of the section ring is best for asymptotics.  This chapter
builds the dictionary carefully and fixes its signs.

\section{Orientation}
\label{sec:ch33-orientation}

\subsection{One degeneration, three languages}
\label{sec:ch33-three-languages}

A test configuration is a flat family over the affine line whose fibre over
$1$ is the given variety and which is equivariant for the multiplicative
group.  That single object is presented in this chapter in three ways, and
each presentation makes a different question easy.

\emph{The total space.}  Writing the degeneration as a scheme
$\pi:\cX\to\mathbf A^1$ makes divisors, blowups, and intersection numbers
available.  Everything numerical in
\Crefrange{chap:bj-energy}{chap:bj-stability}---energies, the
Donaldson--Futaki invariant, the norms $J$ and $I$---is ultimately an
intersection number on a compactification of $\cX$.  The total space is
also where the passage to the trivially valued analytification happens: by
\Cref{chap:trivial-analytification}, the Gauss section realizes
$X^{\An}$ inside $\cX^{\An}$ as the locus of invariant semivaluations
with $v(\varpi)=1$, so that a vertical divisor on $\cX$ becomes a
function on $X^{\An}$.

\emph{The graded algebra.}  Writing the same degeneration as a Rees algebra
turns questions about existence of a model into questions about finite
generation of an algebra.  The flag-ideal presentation
\eqref{eq:flag-ideal} makes the semiring structure
\eqref{eq:flag-semiring} visible, and it is the presentation under which
maxima and sums of the associated functions become sums and products of
ideals.  When we later need to know that a given continuous function is
piecewise linear, this is the language in which the check is finite.

\emph{The filtration.}  Writing the degeneration as a decreasing filtration
of the section ring $R(X,L)$ retains its filtered subspaces and forgets
modifications contracted by a semiample polarization.  The individual
weight multiplicities give a further, coarser invariant.  This is the
language in which asymptotics live: the
Duistermaat--Heckman measure of \Cref{sec:ch33-dh} is a limit of empirical
weight distributions, and it is defined for filtrations far more general
than those coming from a model.  It is also the language that survives the
passage to the complex side in \Cref{chap:rays-filtrations}.

The precise dictionaries involve a polarized model and its closure
convention.  A flag ideal reconstructs its blowup, while a relatively
globally generated vertical divisor on a higher model may contract it.
A finitely generated section algebra reconstructs the ample model of a
semiample polarization;
this is \Cref{thm:ch33-flag-tc-equivalence} and
\Cref{thm:ch33-fg-filtration-tc}.  Filtrations are strictly more general:
\Cref{ex:ch33-harmonic-not-fg} exhibits a multiplicative
linearly bounded filtration that is not of finite type, so that no
semiample test configuration produces it.  A large part of the chapter is
devoted to saying precisely where the equivalence stops.

\subsection{What is proved here and what is deferred}
\label{sec:ch33-scope}

The following are proved in this chapter.

\begin{enumerate}
\item\label{itm:ch33-scope:tc}
 Test configurations form a directed poset under domination, the
 compactification $\overline\cX\to\mathbf P^1$ exists and is projective,
 the extension $\overline\cL$ is unique, and the intersection numbers
 computed on $\overline\cX$ are invariant under domination
 (\Cref{prop:ch33-directed-poset,thm:ch33-compactification,thm:ch33-intersection-invariance}).
\item\label{itm:ch33-scope:flag}
 Flag ideals produce integrally closed blowups of the trivial configuration;
 relatively globally generated vertical divisors give flag ideals in the
 reverse direction, with the precise domination qualifications of
 \Cref{thm:ch33-flag-tc-equivalence}.
\item\label{itm:ch33-scope:norm}
 A relatively semiample test configuration with ample generic polarization
 induces a multiplicative, linearly bounded, finitely generated filtration
 after the integral-closure convention below.
 Conversely a finitely generated filtration produces its saturated
integral-closure class; the exact qualification is stated in
 \Cref{thm:ch33-fg-filtration-tc}.
\item\label{itm:ch33-scope:base}
 Degree-specific base ideals compute the associated flag ideals and the
 function on $X^{\An}$.  After a generating Veronese is fixed, their
 integral closures propagate by powers; the associated function is
 unchanged by integral closure
 (\Cref{prop:ch33-base-ideal-flag,prop:tc-valuative-invariance}).
\item\label{itm:ch33-scope:dh}
 The existence of the Duistermaat--Heckman measure is a cited foundational
 input for every linearly bounded multiplicative filtration of an
 ample section ring.  The resulting measure is a probability measure
 with compact support, is the push-forward of Lebesgue measure on
 an Okounkov body by a concave transform, and transforms in the expected
 way under translation, scaling, and base change
 (\Cref{thm:dh-limit-filtration,thm:ch33-dh-concave-transform,prop:ch33-dh-transformations}).
 For a relatively ample test configuration it is the asymptotic weight distribution of the
 central-fibre representation
 (\Cref{thm:ch33-dh-weight-count}).
\end{enumerate}

The following are \emph{not} proved here.

\begin{enumerate}
\item The identification of $\PL(X)$ with the direct limit of vertical
 Cartier divisors, and of Fubini--Study functions with semiample test
 configurations, belongs to \Cref{chap:bj-fs-pl-psh}; see
 \Cref{thm:tc-pl-correspondence} there.  We construct flag-ideal functions
 here and import the stated injectivity when comparing semiample models.
\item The correspondence with geodesic rays in the space of positively
 curved metrics over $\mathbf C$ is stated in
 \Cref{sec:ch33-geodesic-deferred} and motivated there, but its proof
 requires complex Monge--Amp\`ere theory and is given only in
 \Cref{chap:rays-filtrations}.  Nothing in this chapter depends on it.
\item No asymptotic-slope statement relating a complex energy functional
 along a ray to a non-Archimedean functional is proved or used here.  The
 published sources for such statements, and the exact boundary imposed by
 the 2022 erratum to the Boucksom--Hisamoto--Jonsson energy asymptotics,
 are recorded in \Cref{thm:rf-erratum-audit}; see also
 \Cref{sec:ch33-geodesic-deferred}.
\end{enumerate}

\subsection{Standing hypotheses}
\label{sec:ch33-hypotheses}

Throughout, $X$ is a projective $k$-variety in the sense of
\cite[\S1]{BJ22}, and $L$ is a
$\mathbf Q$-line bundle.  When Hilbert and weight asymptotics are used,
$X$ is integral of dimension $n$ and $L$ is ample.  We write
$R(X,L)=\bigoplus_{m\geq0}R_m$, where
$R_m=H^0(X,mL)$ for sufficiently divisible $m$.

The ground field $k$ is algebraically closed and carries the trivial
absolute value, as in \Cref{chap:trivial-analytification}.  We write
$\varpi$ for the standard coordinate on $\mathbf A^1_k$, so that
$\mathbf G_m=\Spec k[\varpi^{\pm1}]$ and the standard action of
$\mathbf G_m$ on $\mathbf A^1$ is
$t\cdot\varpi=t\varpi$.  When $L$ is ample of dimension $n$, we set
\begin{equation}
\label{eq:ch33-volume}
 V:=(L^n)=\lim_{m\to\infty}\frac{n!}{m^n}\,h^0(X,mL)>0 .
\end{equation}
Fix a denominator $q$ with $qL$ an honest line bundle and retain degrees
$m\in q\mathbf Z_{\geq0}$.  Limits $m\to\infty$ are taken through these
degrees in the usual order, after increasing $q$ when required.
An increasing supremum of degree-normalized functions uses instead the
divisibility order on positive retained degrees; powers make this family
directed.  These are two different indexing conventions, explicitly
distinguished where they are used.

\subsection{Conventions on weights and linearizations}
\label{sec:ch33-weight-conventions}

Two sign choices are made once and used everywhere.

\emph{Weights of a representation.}  If $W$ is a
$k$-vector space with a $\mathbf G_m$-action, then
$W=\bigoplus_{\lambda\in\mathbf Z}W_\lambda$, where $t\in\mathbf G_m$
acts on $W_\lambda$ by $t^\lambda$.  The elements of $W_\lambda$ are
said to have weight $\lambda$.  This is the standard convention and is
forced on us by the identification of a graded algebra with a
$\mathbf G_m$-scheme.

\emph{Weights of a filtration.}  The filtration attached to a test
configuration in \Cref{def:ch33-tc-filtration} is normalized so that
\emph{larger weight means harder to extend}:
$s\in\mathcal F^\lambda R_m$ requires $\varpi^{-\lambda}\bar s$ to be
regular, and $\varpi^{-\lambda}$ has a pole of order $\lambda$ along the
central fibre.  With this normalization the associated norm is
$\lVert s\rVert_\chi=\rme^{-\chi_m(s)}$, consistently with the global rule
$v_K(a)=-\log|a|$ of the Conventions, and twisting the linearization by
the character of weight $c$ sends $\chi_m$ to $\chi_m+mc$ and the
associated function on $X^{\An}$ to itself plus $c$.

\emph{Vertical divisors.}  The sign of $\phi_D$ for a vertical divisor
$D$ is that of the Conventions and of
\eqref{eq:appJ-model-divisor}: an effective vertical divisor gives a
nonnegative function, and $\phi_{\cX_0}=1$.  The two unit tests
\eqref{eq:appJ-special-fibre-test} and the translation test of
\Cref{sec:appJ-rays} are applied to every formula of this chapter at the
point where it is proved, and the worked examples of
\Cref{sec:ch33-examples} check all of them numerically at least once.

\section{Test configurations and compactification}
\label{sec:ch33-test-configurations}

\begin{definition}
\label{def:bj-test-configuration}
A test configuration for $X$ consists of a flat projective morphism
$\pi:\cX\to\mathbf A^1_k$, a $\mathbf G_m$-action on $\cX$ lifting
the standard action on $\mathbf A^1$, and a
$\mathbf G_m$-equivariant isomorphism
\begin{equation}
\label{eq:tc-punctured-trivialization}
 \cX\setminus\cX_0\simeq X\times\mathbf G_m.
\end{equation}
A test configuration $(\cX,\cL)$ for $(X,L)$ additionally carries a
$\mathbf G_m$-linearized $\mathbf Q$-line bundle together with a fixed identification of its restriction over
$1\in\mathbf G_m$ with $L$.
\end{definition}

Some authors package the last datum as an isomorphism
$\cX_1\simeq X$ rather than as
\eqref{eq:tc-punctured-trivialization}.  The two are equivalent: an
isomorphism of the fibre over $1$ propagates over
$\mathbf G_m$ by equivariance, and conversely
\eqref{eq:tc-punctured-trivialization} restricts to an isomorphism over
$1$.  We use the trivialized form because every weight computation below
is a comparison of two extensions of the \emph{same} object over the
punctured line.

\subsection{Elementary structure of the total space}
\label{sec:ch33-tc-structure}

\begin{lemma}\label{lem:ch33-central-fibre-cartier}
Let $\pi:\cX\to\mathbf A^1$ be a test configuration for $X$.  Then:
\begin{enumerate}
\item the central fibre $\cX_0$ is the principal Cartier divisor
 $\{\pi^\ast\varpi=0\}$, and $\pi^\ast\varpi$ is a nonzerodivisor in
 $\cO_{\cX}$;
\item $\cX\setminus\cX_0$ is Zariski dense in $\cX$, and every
 irreducible component of $\cX$ meets it;
\item $\dim\cX=\dim X+1$, and $\cX$ is reduced;
\item for every irreducible closed subvariety $Y\subset X$, the closure
 in $\cX$ of the image of $Y\times\mathbf G_m$ under
 \eqref{eq:tc-punctured-trivialization} is a test configuration for $Y$,
 and this construction induces a bijection between the irreducible
 components of $X$ and those of $\cX$.
\end{enumerate}
\end{lemma}

\begin{proof}
(1)  Flatness of $\pi$ over the regular one-dimensional base
$\mathbf A^1$ means exactly that $\cO_{\cX}$ is a torsion-free
$k[\varpi]$-module, by \Cref{lem:torsion-free-flat-valuation} applied
locally at each point of $\mathbf A^1$; torsion-freeness says that
$\pi^\ast\varpi$ is a nonzerodivisor on every affine chart.  A
nonzerodivisor cuts out an effective Cartier divisor, whose underlying
closed subscheme is $\pi^{-1}(0)=\cX_0$.

(2)  If $Z\subset\cX$ were an irreducible component contained in
$\cX_0$, then $\pi^\ast\varpi$ would belong to the maximal ideal
of the zero-dimensional noetherian local ring $\cO_{\cX,\xi}$ at its
generic point $\xi$.  It would therefore be nilpotent and a zerodivisor,
contradicting (1).  Thus every component meets the open set
$\cX\setminus\cX_0$, which is therefore dense.

(3)  Every component meets the punctured product, and its dimension is
the dimension of its corresponding component of $X$ plus one.
Thus $\dim\cX=\dim X+1$.  For reducedness, a nilpotent local function
becomes zero after inverting $\varpi$, since the punctured product is
reduced.  It is therefore killed by a power of $\varpi$, and flatness
forces it to be zero.  This is also
\cite[Proposition~2.6]{BHJ17} and \cite[\S1.2]{BJ22}.

(4)  Under \eqref{eq:tc-punctured-trivialization} the subscheme
$Y\times\mathbf G_m$ is $\mathbf G_m$-invariant, so its closure
$\cY\subset\cX$ is invariant, and $\cY\to\mathbf A^1$ is projective.
It is flat because $\cY$ is the scheme-theoretic closure of a subscheme
of the generic fibre, hence has no $\varpi$-torsion
(\Cref{lem:flat-closure-generic-fibre}).  Its restriction over
$\mathbf G_m$ is $Y\times\mathbf G_m$ by construction.  Applying this to
the irreducible components $X_\alpha$ of $X$ produces invariant closed
subschemes $\cX_\alpha\subset\cX$ whose union contains the dense open set
$\cX\setminus\cX_0$; since each $\cX_\alpha$ is irreducible of
dimension $\dim X_\alpha+1$ and they are pairwise distinct, they are
exactly the irreducible components of $\cX$.
\end{proof}

\begin{remark}
\label{rem:ch33-flatness-role}
Flatness makes $\cX_0$ a Cartier divisor and makes the total space the
flat closure of its generic fibre.  Both uses are
recorded in \Cref{lem:ch33-central-fibre-cartier}, and neither requires
normality or irreducibility.  Flatness excludes irreducible components of
the total space supported over the origin, but it does not exclude embedded
points of the central fibre.  Their absence is instead the vertical
$S_2$ condition in
\Cref{thm:ch33-integrally-closed-criterion}, hence one half of integral
closedness.
\end{remark}

On an invariant affine chart, equivariance gives every quasi-coherent
invariant sheaf a weight decomposition.  We record the local statement,
together with the gluing hypothesis used repeatedly below.

\begin{lemma}
\label{lem:ch33-weight-decomposition}
Let $A$ be a $k$-algebra with a $\mathbf G_m$-action, that is, a
$\mathbf Z$-grading $A=\bigoplus_{\lambda\in\mathbf Z}A_\lambda$, and
let $M$ be an $A$-module with a compatible $\mathbf G_m$-action.
Then $M=\bigoplus_{\lambda\in\mathbf Z}M_\lambda$ with
$A_\mu M_\lambda\subset M_{\lambda+\mu}$, and a submodule of $M$ is
invariant if and only if it is the direct sum of its weight pieces.  On an
equivariant scheme admitting an invariant affine cover, the same statements
hold for quasi-coherent invariant subsheaves of an equivariant
quasi-coherent sheaf.
\end{lemma}

\begin{proof}
An action of $\mathbf G_m=\Spec k[t^{\pm1}]$ on $M$ is a comodule
structure $\rho:M\to M\otimes_kk[t^{\pm1}]$; writing
$\rho(x)=\sum_\lambda x_\lambda\otimes t^\lambda$ and using the counit and
coassociativity axioms gives $x=\sum_\lambda x_\lambda$ and
$\rho(x_\lambda)=x_\lambda\otimes t^\lambda$, which is the asserted
decomposition; compatibility of $\rho$ with the module structure gives
$A_\mu M_\lambda\subset M_{\lambda+\mu}$.  If $N\subset M$ is invariant
then $\rho(N)\subset N\otimes k[t^{\pm1}]$, so each component
$x_\lambda$ of $x\in N$ lies in $N$; the converse is clear.  Sheafify
on the assumed invariant affine cover.  The lemma asserts no existence
of such a cover for an arbitrary nonnormal scheme.  The action on global
sections, used later, has its weight decomposition directly as a comodule.
\end{proof}

\subsection{Morphisms, domination, and the directed poset}
\label{sec:ch33-domination}

A morphism of test configurations $\mu:\cX'\to\cX$ is a
$\mathbf G_m$-equivariant morphism over $\mathbf A^1$ compatible with
the trivializations \eqref{eq:tc-punctured-trivialization}.  We say that
$\cX'$ \emph{dominates} $\cX$ when such a morphism exists.

\begin{proposition}\label{prop:ch33-directed-poset}
Let $\cX,\cX'$ be test configurations for $X$.
\begin{enumerate}
\item There is at most one morphism $\cX'\to\cX$.
\item If $\cX'$ dominates $\cX$ and $\cX$ dominates $\cX'$, the two
 morphisms are mutually inverse isomorphisms.
\item There is a test configuration $\cX''$ dominating both $\cX$ and
 $\cX'$, obtained as the flat closure in $\cX\times_{\mathbf A^1}\cX'$
 of the diagonally embedded $X\times\mathbf G_m$.
\end{enumerate}
Consequently the isomorphism classes of test configurations for $X$ form
a directed poset.
\end{proposition}

\begin{proof}
(1)  Two morphisms $\mu_1,\mu_2:\cX'\to\cX$ agree on the dense open set
$\cX'\setminus\cX'_0$ by compatibility with the trivializations.  The
locus where they agree is closed in $\cX'$ because $\cX$ is separated,
and it contains a dense subset, hence equals $\cX'$.  Here we use that
$\cX'$ is reduced, by \Cref{lem:ch33-central-fibre-cartier}\,(3), so
that a closed subscheme containing a dense open subscheme is everything.

(2)  By (1), the two composites are the identity, since the identity is one
morphism from each configuration to itself.

(3)  Let $\mathcal Z\subset\cX\times_{\mathbf A^1}\cX'$ be the scheme-theoretic
closure of the image of $X\times\mathbf G_m$ under the diagonal of the two
trivializations.  It is invariant, projective over $\mathbf A^1$, and flat
by \Cref{lem:flat-closure-generic-fibre}, so it is a test configuration; the
two projections are morphisms of test configurations by construction.
Directedness follows, the poset structure being given by (1) and (2).
\end{proof}

\begin{remark}
\label{rem:ch33-domination-rigid}
Assertion (1) makes the formalism rigid: the identification
over $\mathbf G_m$ is part of the data, and it removes all automorphisms.
Without it, the product configuration $X\times\mathbf A^1$ would dominate
itself in as many ways as $X$ has automorphisms, and no directed limit of
vertical divisor groups would be well defined.  This rigidity is used in
\Cref{chap:bj-fs-pl-psh} whenever a PL function is presented on
``a'' determination.
\end{remark}

\subsection{The compactification and its role}
\label{sec:ch33-compactification}

The central fibre $\cX_0$ is the principal Cartier divisor
$(\varpi=0)$.  Gluing \eqref{eq:tc-punctured-trivialization} to
$X\times(\mathbf P^1\setminus\{0\})$ gives a canonical equivariant
compactification
\begin{equation}
\label{eq:tc-compactification}
              \overline\pi:\overline\cX\longrightarrow\mathbf P^1,
\end{equation}
and $\cL$ extends uniquely using the fixed trivialization near infinity.
This compactification is what enters intersection numbers.  Both assertions
require proof, and they are the content of the next two statements.

\begin{theorem}
\label{thm:ch33-compactification}
Let $(\cX,\cL)$ be a test configuration for $(X,L)$.
\begin{enumerate}
\item The scheme $\overline\cX$ obtained by gluing $\cX$ to
 $X\times(\mathbf P^1\setminus\{0\})$ along
 \eqref{eq:tc-punctured-trivialization} is a projective
 $k$-scheme carrying a $\mathbf G_m$-action and a projective
 equivariant morphism $\overline\pi$ to $\mathbf P^1$ restricting to
 $\pi$ over $\mathbf A^1$.
\item There is a unique $\mathbf G_m$-linearized $\mathbf Q$-line bundle
 $\overline\cL$ on $\overline\cX$ restricting to $\cL$ on $\cX$ and
 to the pullback of $L$ on $X\times(\mathbf P^1\setminus\{0\})$, with the fixed compatible
 identifications on the overlap.
\item If $\mu:\cX'\to\cX$ is a morphism of test configurations, it extends
 uniquely to $\overline\mu:\overline{\cX'}\to\overline\cX$, and
 $\overline{\mu^\ast\cL}=\overline\mu^\ast\overline\cL$.
\end{enumerate}
\end{theorem}

\begin{proof}
(1)  The two open pieces $\cX$ and $X\times(\mathbf P^1\setminus\{0\})$
are glued along the common open subscheme $X\times\mathbf G_m$, the
gluing datum being \eqref{eq:tc-punctured-trivialization}; the result is a
scheme $\overline\cX$ with a morphism $\overline\pi$ to
$\mathbf P^1$ obtained by gluing $\pi$ and the second projection, which
agree over $\mathbf G_m$.  Properness is local on the target:
$\overline\pi$ is proper over $\mathbf A^1$ because $\pi$ is, and
proper over $\mathbf P^1\setminus\{0\}$ because $X$ is proper.  A proper
morphism is projective once it has a relatively ample line bundle.
Choose a $\mathbf G_m$-linearized $\pi$-ample bundle $\cM$, as in the
equivariant projectivity input used in \cite[Lemma~1.8]{BJ22}.  Its
punctured-product identification extends it by $p_1^\ast M$, where
$M=\cM|_{\cX_1}$, to a bundle on $\overline\cX$.  This bundle is
relatively ample on both charts, hence relatively ample over
$\mathbf P^1$.  Thus $\overline\pi$ is projective; composing with the
projective base gives projectivity over $k$.
The $\mathbf G_m$-action glues because the gluing datum is equivariant
and the action on $\mathbf P^1\setminus\{0\}$ is the one induced from
$\mathbf A^1$ by $\varpi\mapsto\varpi^{-1}$ on the second chart.

(2)  Existence: glue $\cL$ to $p_1^\ast L$ on
$X\times(\mathbf P^1\setminus\{0\})$ using the isomorphism supplied by
the definition of a test configuration for $(X,L)$.  Uniqueness: if
$\overline\cL_1,\overline\cL_2$ are two such extensions, then
$\overline\cL_1-\overline\cL_2$ is a $\mathbf Q$-line bundle equipped
with a trivialization over each of the two glueing charts; the two
trivializations agree over $X\times\mathbf G_m$ by hypothesis, so they
glue to a global trivialization.  For the linearizations, the same argument
applies to the difference of the two $\mathbf G_m$-linearizations, which
is a character of $\mathbf G_m$ that is trivial over $\mathbf G_m$,
hence trivial.

(3)  Extend $\mu$ by the identity of
$X\times(\mathbf P^1\setminus\{0\})$; the two definitions agree over
$X\times\mathbf G_m$ because $\mu$ is compatible with the
trivializations.  Uniqueness follows as in
\Cref{prop:ch33-directed-poset}\,(1).  The last assertion is the
uniqueness in (2) applied to $\overline\mu^\ast\overline\cL$.
\end{proof}

\begin{warning}{The compactification is not $X\times\mathbf P^1$}
\label{warn:ch33-compactification-not-product}
Only the chart at infinity is a product.  The scheme $\overline\cX$
contains the possibly very singular central fibre $\cX_0$ and, in
general, is not birational to $X\times\mathbf P^1$ by a morphism in the
wrong direction.  What is canonical is the pair
$(\overline\cX,\overline\cL)$ together with the isomorphism over
$\mathbf P^1\setminus\{0\}$; this is what makes the intersection numbers
of \Cref{thm:ch33-intersection-invariance} well defined.
\end{warning}

Why compactify at all?  Because the numerical invariants of the
degeneration are top intersection numbers, and on the noncompact
$\cX$ there is no degree map.  The next theorem is the precise statement
that the compactification does its job.

\begin{theorem}
\label{thm:ch33-intersection-invariance}
Assume $X$ is integral of dimension $n$ and let $(\cX,\cL)$ be a test
configuration for $(X,L)$.
\begin{enumerate}
\item The number $(\overline\cL^{\,n+1})$ is a well-defined rational
 number, computed on the projective $(n+1)$-dimensional scheme
 $\overline\cX$.
\item If $\mu:\cX'\to\cX$ is a morphism of test configurations and
 $\cL'=\mu^\ast\cL$, then
 $({\overline{\cL'}}^{\,n+1})=(\overline\cL^{\,n+1})$.  More generally,
 for $\mathbf Q$-line bundles $\cL_0,\dots,\cL_n$ determined on
 $\cX$,
 \begin{equation}
 \label{eq:ch33-intersection-domination}
  (\overline{\mu^\ast\cL_0}\cdots\overline{\mu^\ast\cL_n})
  =(\overline{\cL_0}\cdots\overline{\cL_n}).
 \end{equation}
\item Twisting the linearization of $\cL$ by the character of weight
 $c\in\mathbf Q$ replaces $\overline\cL$ by
 $\overline\cL+c\,\overline\pi^\ast\{0\}$ and gives
 \begin{equation}
 \label{eq:ch33-intersection-twist}
  \bigl((\overline\cL+c\,\overline\pi^\ast\{0\})^{n+1}\bigr)
  =(\overline\cL^{\,n+1})+(n+1)c\,(L^n).
 \end{equation}
\item If $\cX_d\to\mathbf A^1$ denotes base change along
 $\varpi\mapsto\varpi^d$, followed by integral closure in its generic
 fibre, and $\cL_d$ is the pullback of $\cL$, then
 \begin{equation}
 \label{eq:ch33-intersection-base-change}
  (\overline{\cL_d}^{\,n+1})=d\,(\overline\cL^{\,n+1}).
 \end{equation}
\end{enumerate}
\end{theorem}

\begin{proof}
(1)  By \Cref{thm:ch33-compactification}\,(1), $\overline\cX$ is a
projective $k$-scheme; by
\Cref{lem:ch33-central-fibre-cartier}\,(3) it is integral of dimension
$n+1$.  Top self-intersection of a $\mathbf Q$-line bundle on a
projective scheme is defined by $\mathbf Q$-linear extension of the
Snapper polynomial \cite[Ch.~2]{Ful}, so $(\overline\cL^{\,n+1})$ is a
well-defined rational number.

(2)  The morphism $\overline\mu$ of
\Cref{thm:ch33-compactification}\,(3) is projective, and it is an
isomorphism over $\mathbf P^1\setminus\{0\}$, hence birational.  The
projection formula \Cref{thm:model-projection-formula} gives
\eqref{eq:ch33-intersection-domination}, since
$\overline\mu_\ast[\overline{\cX'}]=[\overline\cX]$.

(3)  A $\mathbf G_m$-linearization of $\cL$ is determined up to the
character group of $\mathbf G_m$, which is $\mathbf Z$; the twist by
weight $c$ changes the glueing of $\cL$ to $p_1^\ast L$ at infinity by
$\varpi^{c}$, hence changes $\overline\cL$ by the divisor
$c\,\overline\pi^\ast\{0\}$, the fibre over $0\in\mathbf P^1$.  Since
$\overline\pi^\ast\{0\}$ is a fibre class, we have
$(\overline\pi^\ast\{0\})^2=0$ and
$(\overline\cL^{\,n}\cdot\overline\pi^\ast\{0\})=(L^n)$ because
$\overline\cL$ restricts to $L$ on the general fibre and intersection
numbers are constant in a flat family.  The binomial expansion of
$(\overline\cL+c\,\overline\pi^\ast\{0\})^{n+1}$ therefore has exactly two
surviving terms, giving \eqref{eq:ch33-intersection-twist}.

(4)  The morphism $\overline{\cX_d}\to\overline\cX$ induced by base change
is finite of degree $d$ onto its image, and
$\overline{\cL_d}$ is its pullback, so
\eqref{eq:ch33-intersection-base-change} is the degree formula for a finite
morphism of degree $d$, again \Cref{thm:model-projection-formula}.
\end{proof}

\begin{remark}
\label{rem:ch33-two-unit-tests}
Formulas \eqref{eq:ch33-intersection-twist} and
\eqref{eq:ch33-intersection-base-change} are the intersection-theoretic
shadow of the translation and dilation rules for Duistermaat--Heckman
measures proved in \Cref{prop:ch33-dh-transformations}.  The consistency of
the three statements is checked explicitly on a nontrivial example in
\Cref{ex:ch33-dnc-p1}; the numerical check there is
$(\overline\cL^{2})=-c^2$ against a first moment $-c^2/2$, in agreement
with \eqref{eq:cma-dh-first-intersection}.
\end{remark}

A morphism of test configurations is equivariant over $\mathbf A^1$ and
the identity on the punctured product.  Two test configurations for
$(X,L)$ are \emph{equivalent} if their polarizations agree after pulling
back to a third test configuration dominating both; this is
\cite[Definition~6.1]{BHJ17}, and it is the equivalence relation under
which the associated vertical-divisor function on $X^{\An}$ is a
complete invariant \cite[Theorem~2.31(i)]{BJ22}.  For a semiample
polarization this agrees with the section-supremum function, by
\Cref{prop:tc-valuative-invariance}.  Pulling back by
$\varpi\mapsto\varpi^d$, and then taking integral closure, is the
degree-$d$ base change.  The raw base change rescales weights by $d$;
subsequent integral closure may enlarge individual filtered pieces
(\Cref{prop:ch33-transformations}).

\subsection{Integrally closed test configurations}
\label{sec:ch33-integral-closure}

\begin{definition}
\label{def:integrally-closed-tc}
A test configuration $\cX$ is \emph{integrally closed} if it is integrally
closed in its generic fibre $X_{k(\varpi)}$.  Equivalently, on each affine
chart $\Spec A$, the $k[\varpi]$-algebra $A$ is integrally closed in
$A\otimes_{k[\varpi]}k(\varpi)$.
\end{definition}

This condition is weaker than normality when $X$ is not normal.  The
integral closure of any test configuration in its generic fibre is finite
and is characterized as the unique integrally closed finite domination.
Thus replacing a test configuration by its integral closure loses no
valuation-theoretic information.  We now prove the three facts about
integral closure that are used below.

\begin{proposition}\label{prop:ch33-integral-closure-exists}
Let $\cX$ be a test configuration for $X$.
\begin{enumerate}
\item The integral closure of $\cO_{\cX}$ inside the sheaf of regular
 functions on the generic fibre defines a finite,
 $\mathbf G_m$-equivariant morphism
 $\widetilde\cX\to\cX$ which is an isomorphism over
 $\mathbf G_m$; the scheme $\widetilde\cX$ is an integrally closed test
 configuration for $X$.
\item $\widetilde\cX$ is the unique integrally closed test configuration
 dominating $\cX$ by a finite morphism.
\item Integrally closed test configurations are cofinal in the directed
 poset of \Cref{prop:ch33-directed-poset}.
\item If $\cX$ is integrally closed and $\mu:\cX'\to\cX$ is any
 morphism of test configurations, then $\mu_\ast\cO_{\cX'}=\cO_{\cX}$.
\end{enumerate}
\end{proposition}

\begin{proof}
(1)  The scheme $\cX$ is of finite type over a field, hence Nagata and
excellent, so the integral closure of $\cO_{\cX}$ inside the generic-fibre structure
sheaf is a coherent $\cO_{\cX}$-algebra, by the
finiteness statement recorded in \Cref{def:appA-normalized-blowup} and
\cite[\href{https://stacks.math.columbia.edu/tag/035S}{Tag 035S}]
{stacks-project}.  Its relative $\Spec$ is finite over $\cX$, hence
projective over $\mathbf A^1$; it is flat over $\mathbf A^1$ because it
is a subsheaf of a sheaf of $k(\varpi)$-algebras, hence torsion free.  The
construction is canonical, so it inherits the $\mathbf G_m$-action, and
it is an isomorphism over $\mathbf G_m$ because $\cX$ restricted there
is $X\times\mathbf G_m$.  The relative integral-closure construction and its unchanged
punctured product are the precise input of
\cite[Definition~1.10 and the paragraph following Example~1.11]{BJ22}.
They concern closure inside generic regular functions, and do not
normalize a nonnormal $X$.

(2)  If $\cX'\to\cX$ is finite with $\cX'$ integrally closed and an
isomorphism on generic fibres, then $\cO_{\cX'}$ is a finite
$\cO_{\cX}$-algebra inside the generic-fibre function sheaf and is
integrally closed; hence it contains the integral closure and is contained
in it.

(3)  Given $\cX$, the configuration $\widetilde\cX$ dominates it and is
integrally closed.  Directedness then makes the integrally closed ones
cofinal.

(4)  This is Zariski's main theorem in the form used by
\cite[Lemma~1.12]{BJ22}, and the proof is short enough to give: by
properness of $\mu$ and coherence of higher direct images,
$\mu_\ast\cO_{\cX'}$ is a coherent $\cO_{\cX}$-algebra.  A section over
an open $U\subset\cX$ is a rational function on $U$, regular on the
generic fibre, and integral over $\cO_U$; integral closedness of $\cX$
forces it to lie in $\cO_U$.
\end{proof}

\refstepcounter{theorem}\label{thm:ch33-integrally-closed-criterion}
\begin{quote}
\textbf{Foundational input \thetheorem\ (vertical Serre criterion; source
boundary).}
A test configuration $\cX$ is integrally closed if and only if it is
\emph{vertically $R_1$}, meaning that $\cX$ is regular at each generic
point of $\cX_0$, and \emph{vertically $S_2$}, meaning that
$\operatorname{depth}\cO_{\cX,\xi}\geq\min\{2,\dim\cO_{\cX,\xi}\}$ for
every $\xi\in\cX_0$.  Moreover $\cX$ is vertically $S_2$ if and only
if $\cX_0$ has no embedded points, and $\cX_0$ is reduced if and only if
$\cX$ is integrally closed with generically reduced central fibre.
These statements are imported from \cite[Theorem~1.13, Lemma~1.15
and Corollary~1.16]{BJ22}; no proof is
claimed here.  The statement is used only through its first consequence, that
the local ring of $\cX$ at a generic point of $\cX_0$ is a discrete
valuation ring.  That consequence is what attaches divisorial valuations to
components of the central fibre.
\end{quote}

\begin{definition}
\label{def:ch33-component-valuation}
Let $\cX$ be an integrally closed test configuration for an integral
$X$, and let $E$ be an irreducible component of $\cX_0$.  By
\Cref{thm:ch33-integrally-closed-criterion} the local ring of $\cX$ at
the generic point of $E$ is a discrete valuation ring, defining
$\ord_E$ on $k(\cX)^\times$.  Set
\begin{equation}
\label{eq:ch33-bE}
 b_E:=\ord_E(\varpi)=\ord_E(\cX_0)\in\mathbf Z_{>0},
 \qquad
 w_E:=b_E^{-1}\ord_E ,
\end{equation}
and let $v_E\in X^{\divis}$ be the restriction of $w_E$ to
$k(X)\subset k(X)(\varpi)=k(\cX)$.  This may be the trivial valuation,
which belongs to $X^{\divis}$ by the convention of
\Cref{def:ch32-divisorial}.
\end{definition}

\begin{lemma}\label{lem:ch33-gauss-component}
With the notation of \Cref{def:ch33-component-valuation}, the Gauss
extension \eqref{eq:gauss-extension-triv} satisfies
$\sigma(v_E)=w_E$.  Conversely, a valuation $v$ on $X$ is divisorial
if and only if $v=v_E$ for some irreducible component $E$ of the central
fibre of some integrally closed test configuration.
\end{lemma}

\begin{proof}
The valuation $w_E$ is invariant under the $\mathbf G_m(k)$-action, since
$E$ is invariant, and $w_E(\varpi)=1$ by construction.  By
\Cref{chap:trivial-analytification}, the Gauss extension is characterized
as the unique invariant extension of $v$ with value $1$ on $\varpi$;
this is \cite[Lemma~4.2]{BHJ17} as quoted in \cite[\S1.3]{BJ22}.  Both
$w_E$ and $\sigma(v_E)$ restrict to $v_E$ on $k(X)$, are invariant,
and take the value $1$ at $\varpi$; hence they agree.  The converse
statement is \cite[Theorem~4.6]{BHJ17}, quoted in
\cite[Lemma~1.20]{BJ22}; it rests on Zariski's theorem that a divisorial
valuation is realized on a projective birational model.
\end{proof}

\begin{lemma}
\label{lem:ch33-effectivity}
Let $\cX$ be an integrally closed test configuration and let
$\mathfrak a$ be a vertical fractional ideal on $\cX$, that is, a
coherent $\mathbf G_m$-invariant fractional ideal sheaf that is trivial
outside $\cX_0$.  Then
$\mathfrak a\subset\cO_{\cX}$ if and only if $\ord_E(\mathfrak a)\geq0$
for every irreducible component $E$ of $\cX_0$.  Consequently a vertical
$\mathbf Q$-Cartier divisor $D$ on $\cX$ is effective if and only if
$\phi_D(v_E)\geq0$ for every such $E$.
\end{lemma}

\begin{proof}
Necessity is clear.  For sufficiency, let
$\mathfrak p=\{f\in\cO_{\cX}\mid f\mathfrak a\subset\cO_{\cX}\}$ be the
ideal of poles and let $P=V(\mathfrak p)$ be the polar scheme, which is
supported in $\cX_0$.  If $P\ne\varnothing$, pick an associated point
$\xi$ of $P$.  The needed fact, proved in \cite[Lemma~1.19]{BJ22} by the
determinant trick, is that for an integrally closed $\cX$ every associated
point of a polar scheme is a generic point of $\cX_0$; we reproduce the
argument.  By definition of an associated point there is
$f\in\cO_{\cX,\xi}$ with $f\notin\mathfrak p_\xi$ and
$f\mathfrak m_\xi\subset\mathfrak p_\xi$.  Then
$f\mathfrak m_\xi\mathfrak a_\xi$ is an ideal of $\cO_{\cX,\xi}$, so it
is contained in $\mathfrak m_\xi$ or equal to $\cO_{\cX,\xi}$.  In the
first case each element of $f\mathfrak a_\xi$ preserves the finite
module $\mathfrak m_\xi$.  This module is faithful, since it contains
the nonzerodivisor $\varpi$.  The determinant trick therefore makes
each such element integral over $\cO_{\cX,\xi}$, hence contained in it by integral closedness,
contradicting $f\notin\mathfrak p_\xi$.  In the second case
$\mathfrak m_\xi$ is invertible and $\cO_{\cX,\xi}$ is a discrete
valuation ring, so $\xi$ is a generic point of $\cX_0$, that is,
$\xi$ is the generic point of some component $E$.  But then
$\ord_E(\mathfrak a)\geq0$ gives
$\mathfrak a_\xi\subset\cO_{\cX,\xi}$, so $\xi\notin P$, a
contradiction.  Hence $P=\varnothing$ and
$\mathfrak a\subset\cO_{\cX}$.

For the last assertion, pass to a multiple so that $D$ is Cartier and
apply the first part to $\mathfrak a=\cO_{\cX}(-D)$, using
$\sigma(v_E)=w_E$ from \Cref{lem:ch33-gauss-component} and
\[
 \phi_D(v_E)=\sigma(v_E)(D)=b_E^{-1}\ord_E(D).
\]
The two occurrences of $\ord_E$ here are of different kinds and must not
be conflated: $\ord_E(\mathfrak a)$ is the order of the \emph{ideal}
$\mathfrak a$ along $E$, while $\ord_E(D)$ is the coefficient of $E$
in the \emph{divisor} $D$.  For $\mathfrak a=\cO_{\cX}(-D)$ the two
agree, $\ord_E(\mathfrak a)=\ord_E(D)=b_E\phi_D(v_E)$, so the first part
gives $\cO_{\cX}(-D)\subseteq\cO_{\cX}$, that is $D\geq0$, exactly when
$\phi_D(v_E)\geq0$ for every $E$.  This is the sign convention
\eqref{eq:mp-divisor-potential} used throughout the book.
\end{proof}

\begin{remark}
\label{rem:ch33-normalization-vs-integral-closure}
If $X$ is normal, then a test configuration is integrally closed if and
only if it is normal, and the integral closure coincides with the
normalization.  In general the normalization $\cX^\nu$ is a test
configuration for the normalization $X^\nu$, not for $X$; the morphism
$\cX^\nu\to\cX$ factors through $\widetilde\cX$, and
$\cX^\nu\to\widetilde\cX$ is an isomorphism over each generic point of
$\widetilde\cX_0$.  Consequently the components of $\widetilde\cX_0$ and
of $\cX^\nu_0$ define the same set of divisorial valuations
(\cite[Remark~1.22]{BJ22}), which is why no statement below is sensitive to
the choice.
\end{remark}

\begin{example}
\label{ex:product-dnc-test-configurations}
The product $X\times\mathbf A^1$, with trivial action on $X$, is the
trivial test configuration.  If $\mathfrak a\subset\cO_X$ defines a closed
subscheme $Z$, the blowup of
$\mathfrak a+(\varpi)\subset\cO_{X\times\mathbf A^1}$ is the deformation to
the normal cone of $Z$.  Its central fibre contains the strict transform of
$X\times\{0\}$ and the exceptional divisor.  If $L$ is ample, the pullback of $L$ minus a sufficiently small
positive rational multiple of the exceptional divisor gives a relatively
ample polarization.

We record the standard supplement, which we use in
\Cref{ex:ch33-dnc-p1}.  The exceptional divisor of this blowup is the
projective completion $\mathbf P(C_{Z/X}\oplus1)$ of the normal cone
$C_{Z/X}=\Spec_Z\bigoplus_{q\geq0}\mathfrak a^q/\mathfrak a^{q+1}$, the
strict transform of $X\times\{0\}$ is the blowup of $Z$ in $X$, and
the central fibre is the sum of these two effective Cartier divisors
\cite[\S1.2]{BJ22}.  When $X$ is smooth and $Z$ is a smooth point of a
smooth curve, the strict transform is $X$ and the exceptional component
is $\mathbf P^1$, meeting transversally.  Both are $\mathbf P^1$ when
$X=\mathbf P^1$.
\end{example}

\begin{example}
\label{ex:ch33-product-nontrivial-action}
Let $\rho:\mathbf G_m\to\Aut(X,L)$ be a one-parameter subgroup, and let
$\cX=X\times\mathbf A^1$ with the action
$t\cdot(x,\varpi)=(\rho(t)x,t\varpi)$, and
$\cL=p_1^\ast L$ with the linearization induced by $\rho$.  The
underlying scheme is a product, but the test configuration is not the
trivial one: the trivialization
\eqref{eq:tc-punctured-trivialization} is the equivariant one, which
differs from the naive product trivialization by $\rho$.  This is the
first place where forgetting the trivialization would collapse a nontrivial
degeneration to a trivial one; the resulting filtration and
Duistermaat--Heckman measure are computed in \Cref{ex:p1-product-dh}.
\end{example}

\section{Flag ideals and Rees algebras}
\label{sec:ch33-flag-ideals}

We now replace the total space by an ideal.  The gain is that ideals form a
semiring, that their behaviour under valuations is completely understood
(\Cref{sec:appA-integral-dependence}), and that finite generation of a Rees
algebra is a checkable condition.  The loss is that only test configurations
dominating the trivial one are visible; by
\Cref{prop:ch33-directed-poset} this loses no information, since every test
configuration is dominated by one that dominates the trivial one.

\subsection{Vertical fractional ideals}
\label{sec:ch33-vertical-ideals}

\begin{definition}
\label{def:bj-flag-ideal}
A flag ideal on $X\times\mathbf A^1$ is a coherent,
$\mathbf G_m$-invariant fractional ideal, trivial over
$\mathbf G_m$.  After multiplying by a power of $\varpi$, it has a
finite expression
\begin{equation}
\label{eq:flag-ideal}
 \mathfrak a=\mathfrak a_0+
 \mathfrak a_1\varpi+\cdots+\mathfrak a_N\varpi^N,
 \qquad
 \mathfrak a_0\subset\mathfrak a_1\subset\cdots\subset
 \mathfrak a_N=\cO_X .
\end{equation}
\end{definition}

The finite expression is not an extra hypothesis but a consequence of the
definition; we prove it, since the argument fixes the indexing used in every
computation below.

\begin{lemma}
\label{lem:ch33-flag-weight-decomposition}
Let $\mathfrak a$ be a coherent $\mathbf G_m$-invariant fractional ideal
on $X_{\triv}:=X\times\mathbf A^1$, trivial over $\mathbf G_m$.  Then
there is a unique decreasing family of coherent fractional ideals
$\mathfrak a_\lambda\subset\cO_X$, indexed by $\lambda\in\mathbf Z$,
such that
\begin{equation}
\label{eq:ch33-flag-weight}
 \mathfrak a=\sum_{\lambda\in\mathbf Z}\mathfrak a_\lambda\varpi^{-\lambda},
 \qquad
 \mathfrak a_\lambda=\cO_X\ (\lambda\ll0),
 \qquad
 \mathfrak a_\lambda=0\ (\lambda\gg0).
\end{equation}
For every $v\in X^{\An}$ with Gauss extension $\sigma(v)$ as in
\eqref{eq:gauss-extension-triv},
\begin{equation}
\label{eq:ch33-sigma-flag}
 \sigma(v)(\mathfrak a)=\min_{\lambda}\{v(\mathfrak a_\lambda)-\lambda\}.
\end{equation}
\end{lemma}

\begin{proof}
Work on an affine chart $\Spec A\subset X$, so that
$\Spec A[\varpi]\subset X_{\triv}$, and let $I\subset A[\varpi,\varpi^{-1}]$
be the corresponding fractional ideal.  The $\mathbf G_m$-action grades
$A[\varpi^{\pm1}]$ with $A$ in weight $0$ and $\varpi$ in weight
$-1$, as fixed in \Cref{rem:ch33-weight-sign}.  By
\Cref{lem:ch33-weight-decomposition}, $I$ is the direct sum of its weight
pieces, and the piece of weight $\lambda$ has the form
$\mathfrak a_\lambda\varpi^{-\lambda}$ for a unique $A$-submodule
$\mathfrak a_\lambda\subset A$; it is an $A$-module
because $A$ has weight $0$.  Multiplication by $\varpi$ maps
$I$ into $I$, and it shifts weights by one; concretely,
$\varpi\cdot\mathfrak a_\lambda\varpi^{-\lambda}
=\mathfrak a_\lambda\varpi^{-(\lambda-1)}\subset I$ gives
$\mathfrak a_\lambda\subset\mathfrak a_{\lambda-1}$, so the family is
decreasing in $\lambda$.  Coherence of $\mathfrak a$ forces finitely
many generators, whose weights are bounded, so
$\mathfrak a_\lambda=0$ for $\lambda\gg0$; triviality over
$\mathbf G_m$ means that $I$ generates the unit ideal after inverting
$\varpi$, which combined with the decreasing property gives
$\mathfrak a_\lambda=\cO_X$ for $\lambda\ll0$.  The decomposition is
independent of the chart because it is the weight decomposition, so the
$\mathfrak a_\lambda$ glue.  Setting $j=-\lambda$ and multiplying by
$\varpi^{N}$ for $N$ large yields \eqref{eq:flag-ideal}.

For \eqref{eq:ch33-sigma-flag}, a local generating set of $\mathfrak a$ at
the center of $\sigma(v)$ may be taken to consist of elements
$f\varpi^{-\lambda}$ with $f$ in a local generating set of
$\mathfrak a_\lambda$, by the weight decomposition.  Then
$\sigma(v)(f\varpi^{-\lambda})=v(f)-\lambda$ by
\eqref{eq:gauss-extension-triv}, and the value of the ideal is the minimum
over generators, by \eqref{eq:ideal-evaluation-triv}.
\end{proof}

\begin{remark}
\label{rem:ch33-two-indexings}
Formula \eqref{eq:flag-ideal} indexes by $j\geq0$ with
$\mathfrak a_N=\cO_X$, while \eqref{eq:ch33-flag-weight} indexes by
$\lambda\in\mathbf Z$ with $\mathfrak a_\lambda=\cO_X$ for
$\lambda\ll0$.  They are related by $j=-\lambda$ and a global shift by a
power of $\varpi$, which shifts the associated function by a constant.  We
use whichever is more convenient and always display the convention in
force.
\end{remark}

\subsection{The function of a flag ideal}
\label{sec:ch33-flag-function}

The blowup of a flag ideal, followed by integral closure in its unchanged
generic fibre, is an integrally closed test configuration dominating the
trivial one.  When $X$ is normal this is the usual normalized blowup.
Conversely, an integrally closed test configuration that dominates the
trivial one, together with a relatively semiample vertical divisor on it,
gives the same function as a flag ideal, and dominates the resulting
blowup after clearing denominators; this is the content of
\Cref{thm:ch33-flag-tc-equivalence}.  For normal integral $X$, the passage from
the ideal to the model runs through the normalized Rees algebra
\begin{equation}
\label{eq:flag-rees-algebra}
   \mathcal R(\mathfrak a)
      :=\bigoplus_{q\geq0}\overline{\mathfrak a^q}\,u^q ,
\end{equation}
whose relative $\Proj$ over $X_{\triv}$ is the normalized blowup, by
\eqref{eq:appA-normalized-rees}.  The bar is essential: valuations see
integral closure, since
\begin{equation}
\label{eq:valuative-integral-closure}
 \overline{\mathfrak a}
   =\{f\mid w(f)\geq w(\mathfrak a)
             \text{ for every centered divisorial }w\}.
\end{equation}
In this normal integral paragraph, $w$ ranges over the divisorial
valuations of $k(X)(\varpi)$ centered on $X\times\mathbf A^1$, not merely
over valuations of $k(X)$.  For nonnormal $X$, use the integral closure
of the blowup inside its unchanged generic fibre instead of the full
normalization of its Rees algebra.
Equation \eqref{eq:valuative-integral-closure} is
\Cref{thm:appA-valuative-integral-closure}, applied on the excellent scheme
$X_{\triv}$.

For the sign convention of this book, define the flag-ideal function by
\begin{equation}
\label{eq:flag-ideal-function}
 \phi_{\mathfrak a}(v)
   :=-\sigma(v)(\mathfrak a)
   =\max_{0\leq j\leq N}
        \{\log|\mathfrak a_j|(v)-j\},
\end{equation}
where $\log|\mathfrak b|(v)=-v(\mathfrak b)$.  If one writes the weight
decomposition with $\varpi^{-\lambda}$, the last term becomes
$\log|\mathfrak a_\lambda|+\lambda$.  Formula
\eqref{eq:flag-ideal-function}, rather than an isolated sign, is our
definition.

Three unit tests confirm the sign.  For $\mathfrak a=\cO_X$ one gets
$\phi_{\mathfrak a}=0$.  For $\mathfrak a=(\varpi)$, the weight
decomposition has the single term $\mathfrak a_1=\cO_X$, giving
$\phi_{\mathfrak a}=-1$; since the blowup of $(\varpi)$ is the trivial
configuration with the vertical divisor $-\cX_0$, this agrees with
$\phi_{\cX_0}=1$ of \eqref{eq:appJ-special-fibre-test}.  For
$\mathfrak a=\mathfrak b+(\varpi)$ with $\mathfrak b\subset\cO_X$ an
honest ideal, one gets
$\phi_{\mathfrak a}=\max\{\log|\mathfrak b|,-1\}$, a bounded continuous
function which is $0$ off $V(\mathfrak b)$; this is the standard
truncation used in the density argument of
\Cref{thm:triv-divisorial-density}.

\begin{lemma}\label{lem:ch33-flag-continuity}
For every flag ideal $\mathfrak a$, the function $\phi_{\mathfrak a}$ is
a bounded continuous real-valued function on $X^{\An}$, with the following
scaling formula for \eqref{eq:valuation-scaling}:
\begin{equation}
\label{eq:ch33-flag-scaling}
 \phi_{\mathfrak a}(tv)
 =\max_\lambda\{t\log|\mathfrak a_\lambda|(v)+\lambda\},
 \qquad t>0 ,
\end{equation}
in the indexing of \eqref{eq:ch33-flag-weight}.
\end{lemma}

\begin{proof}
First multiply the fractional flag ideal by a power of $\varpi$ to
obtain the nonnegative indexing \eqref{eq:flag-ideal}.  This adds a
constant to its function, so it suffices to prove continuity and
boundedness in that indexing.
Each $v\mapsto\log|\mathfrak a_\lambda|(v)=-v(\mathfrak a_\lambda)$ is
continuous with values in $[-\infty,0]$ by
\eqref{eq:ideal-tropical-identities}, and the maximum in
\eqref{eq:flag-ideal-function} is finite because the top term
$\mathfrak a_N=\cO_X$ contributes the finite value $-N$; the maximum is
bounded above by $0$ and below by $-N$.  A finite maximum of continuous
extended-real functions with one finite member is continuous and finite.
Formula \eqref{eq:ch33-flag-scaling} is immediate from
$(tv)(\mathfrak a_\lambda)=t\,v(\mathfrak a_\lambda)$.
\end{proof}

For a flag-ideal model dominating the product, the geometric and algebraic
constructions fit together as in \cref{fig:tc-rees-geometry}.

\begin{figure}[htbp]
\centering
\begin{tikzpicture}[x=1cm,y=1cm,>=Stealth,
  box/.style={draw,rounded corners=3pt,align=center,inner sep=5pt,font=\small}]
  \node[box] (flag) at (-2.75,3.45)
    {$\mathfrak a=\sum_{j=0}^{N}\mathfrak a_j\varpi^j$\\flag ideal};
  \node[box] (rees) at (2.75,3.45)
    {$\mathcal R(\mathfrak a)
      =\displaystyle\bigoplus_{q\geq0}\overline{\mathfrak a^q}\,u^q$\\
      normalized Rees algebra};
  \node[box] (tc) at (2.75,1.75)
    {$\mathcal X_{\mathfrak a}
      =\operatorname{Proj}_{X\times\mathbf A^1}\mathcal R(\mathfrak a)$\\
      test configuration};
  \node[box] (product) at (2.75,0) {$X\times\mathbf A^1$};
  \node[box] (base) at (-2.75,0) {$\mathbf A^1$};

  \draw[-{Stealth[length=2.2mm]},thick] (flag)--(rees)
    node[midway,above,font=\scriptsize] {powers and integral closure};
  \draw[-{Stealth[length=2.2mm]},thick] (rees)--(tc)
    node[midway,right,font=\scriptsize] {relative $\operatorname{Proj}$};
  \draw[-{Stealth[length=2.2mm]},thick] (tc)--(product)
    node[midway,right,font=\small] {$\nu$};
  \draw[-{Stealth[length=2.2mm]},thick] (product)--(base)
    node[midway,above,font=\small] {$\operatorname{pr}_2$};

  \draw[decorate,decoration={brace,amplitude=4pt},thick]
    (1.5,-0.48)--(4.0,-0.48)
    node[midway,below=5pt,font=\scriptsize,align=center]
      {$\nu$ is an isomorphism over $\mathbf G_m$;\\
       exceptional components lie over $0$};
\end{tikzpicture}
\caption{The normalized Rees construction for a flag ideal.  Integral
closure leaves the punctured product unchanged for normal $X$;
for nonnormal $X$ use relative integral closure in the generic fibre.
Relative
$\operatorname{Proj}$ records the vertical modification over $0$.  The
diagram is scoped to the flag-ideal test configurations discussed in this
section.}
\label{fig:tc-rees-geometry}
\end{figure}

\begin{lemma}\label{lem:flag-semiring-identities}
For flag ideals $\mathfrak a,\mathfrak b$,
\begin{equation}
\label{eq:flag-semiring}
 \phi_{\mathfrak a\mathfrak b}
   =\phi_{\mathfrak a}+\phi_{\mathfrak b},
 \qquad
 \phi_{\mathfrak a+\mathfrak b}
   =\max\{\phi_{\mathfrak a},\phi_{\mathfrak b}\}.
\end{equation}
The associated function is unchanged by integral closure.
\end{lemma}

\begin{proof}
Gauss extension is a valuation.  Hence it sends products of ideals to sums
and sums of ideals to minima.  Multiplication by $-1$ turns the latter
minimum into a maximum.  Equality after integral closure follows from
\eqref{eq:valuative-integral-closure}.

In detail: by \Cref{prop:appA-ideal-arithmetic}, applied to the valuation
$\sigma(v)$ on $X_{\triv}$, one has
$\sigma(v)(\mathfrak a\mathfrak b)
=\sigma(v)(\mathfrak a)+\sigma(v)(\mathfrak b)$ and
$\sigma(v)(\mathfrak a+\mathfrak b)
=\min\{\sigma(v)(\mathfrak a),\sigma(v)(\mathfrak b)\}$; negating gives
\eqref{eq:flag-semiring}.  For the last assertion, integral closure does not
change the value of any divisorial valuation centered on $X_{\triv}$, by
\Cref{cor:appA-integral-closure-values}, and every Gauss extension
$\sigma(v)$ with $v\in X^{\divis}$ is such a valuation, by
\Cref{lem:ch33-gauss-component}; two continuous functions agreeing on the
dense set $X^{\divis}$ of \Cref{thm:triv-divisorial-density} agree
everywhere.
\end{proof}

\begin{proposition}
\label{prop:ch33-flag-faithful}
Let $\mathfrak a,\mathfrak b$ be flag ideals on $X_{\triv}$ with
$X$ integral.  Then
\begin{equation}
\label{eq:ch33-flag-faithful}
 \phi_{\mathfrak a}\leq\phi_{\mathfrak b}
 \iff \mathfrak a\subset\overline{\mathfrak b},
 \qquad
 \phi_{\mathfrak a}=\phi_{\mathfrak b}
 \iff \overline{\mathfrak a}=\overline{\mathfrak b}.
\end{equation}
In particular $\mathfrak a\mapsto\phi_{\mathfrak a}$ is injective on
integrally closed flag ideals.
\end{proposition}

\begin{proof}
By \eqref{eq:flag-ideal-function}, $\phi_{\mathfrak a}\leq\phi_{\mathfrak b}$
on $X^{\An}$ is the same as
$\sigma(v)(\mathfrak a)\geq\sigma(v)(\mathfrak b)$ for all
$v\in X^{\An}$.  To apply the divisorial valuative criterion on
$X_{\triv}$, let $w$ be any divisorial valuation centered there.  If
$w(\varpi)=0$, both flag ideals have value zero, since they become trivial
after inverting $\varpi$.  If $w(\varpi)>0$, rescale to
$w(\varpi)=1$ and put $v=w|_{k(X)}\in X^{\An}$.  Writing a flag ideal in
weight form gives directly
\[
 w(\mathfrak a)=\min_\lambda
   \{v(\mathfrak a_\lambda)-\lambda\}
 =\sigma(v)(\mathfrak a),
\]
and similarly for $\mathfrak b$.  Thus the assumed inequality on all
$v\in X^{\An}$ yields the required inequality for every divisorial $w$,
without any orbit-closure assertion.  Now
\eqref{eq:valuative-integral-closure} translates the resulting inequality
into $\mathfrak a\subset\overline{\mathfrak b}$.  The second equivalence
follows by symmetry, and injectivity is the case
$\mathfrak a=\overline{\mathfrak a}$,
$\mathfrak b=\overline{\mathfrak b}$.
\end{proof}

\subsection{The dictionary in both directions}
\label{sec:ch33-flag-dictionary}

\begin{theorem}
\label{thm:ch33-flag-tc-equivalence}
Let $X$ be a projective $k$-variety.
\begin{enumerate}
\item\label{itm:ch33-flag-tc:blowup}
 Let $\mathfrak a$ be a flag ideal on $X_{\triv}$.  The blowup
 $\operatorname{Bl}_{\mathfrak a}X_{\triv}$, followed by integral closure
 in its generic fibre, is an integrally closed test configuration
 $\cX_{\mathfrak a}$ dominating $X_{\triv}$, and
 $\mathfrak a\cO_{\cX_{\mathfrak a}}=\cO_{\cX_{\mathfrak a}}(D)$ for a
 vertical Cartier divisor $D$ that is relatively ample over
 $X_{\triv}$.  Moreover $\phi_D=\phi_{\mathfrak a}$.
\item\label{itm:ch33-flag-tc:pushforward}
 Conversely, let $\mu:\cX\to X_{\triv}$ be a test configuration
 dominating the trivial one, and let $D$ be a vertical Cartier divisor on
 $\cX$ with $\cO_{\cX}(D)$ $\mu$-globally generated.  Then
 $\mathfrak a:=\mu_\ast\cO_{\cX}(D)$ is a flag ideal,
 $\mathfrak a\cO_{\cX}=\cO_{\cX}(D)$, and
 $\phi_{\mathfrak a}=\phi_D$.
\item\label{itm:ch33-flag-tc:bijection}
 The two constructions agree at the following precise level: the
 flag ideal produced in (2) from the
 datum produced in (1) has the same integral
 closure as $\mathfrak a$, and the pair
 $(\cX_{\mathfrak a},D)$ produced in
 (1) from the datum produced in
 (2) is dominated by the integral closure of $\cX$,
 compatibly with the pullback of $D$.
\end{enumerate}
\end{theorem}

\begin{proof}
(1)
By \Cref{prop:appA-blowup-proj}, the blowup of $\mathfrak a$ is the
relative $\Proj$ of $\bigoplus_q\mathfrak a^q$, it is projective over
$X_{\triv}$, and the pullback ideal $\mathfrak a\cO$ is invertible,
say equal to $\cO(D)$.  With this convention $D$ is anti-effective for
an ordinary ideal, up to the global $\varpi$-shift of
\Cref{rem:ch33-two-indexings}.  The divisor
$D$ is vertical because $\mathfrak a$ is trivial over
$\mathbf G_m$, and it is invariant because $\mathfrak a$ is.  Relative
ampleness of $D$, equivalently of $\cO(1)$ for the $\Proj$
presentation, is the standard property of a blowup.  Taking integral
closure in the generic fibre changes neither the generic fibre nor the
invertibility of $\mathfrak a\cO$, by
\Cref{prop:ch33-integral-closure-exists}\,(1), and produces an integrally
closed test configuration.  Finally, for any $v\in X^{\An}$, the Gauss
extension $\sigma(v)$ has a center on $\cX_{\mathfrak a}$ by
properness, and $\sigma(v)(\mathfrak a)=\sigma(v)(\mathfrak a\cO)
=\sigma(v)(\cO(D))$, so
$\phi_{\mathfrak a}=-\sigma(v)(\cO(D))=\phi_D$ by
\eqref{eq:flag-ideal-function} and the definition
of $\phi_D$ recalled in \Cref{sec:ch33-weight-conventions}.  The sign is
fixed by the unit test $\mathfrak a=(\varpi)$: there $D=-\cX_0$,
$\cO(D)=(\varpi)$, $\sigma(v)(\cO(D))=1$ and
$\phi_{\mathfrak a}=-1=\phi_{-\cX_0}$, whereas $\cO(-D)=(\varpi^{-1})$
would give $+1$.

(2)
The sheaf $\mu_\ast\cO_{\cX}(D)$ is coherent because $\mu$ is
projective, invariant because everything in sight is, and trivial over
$\mathbf G_m$ because $\mu$ is an isomorphism there and $D$ is
vertical.  It is a fractional ideal because $\cO_{\cX}(D)$ is a
subsheaf of the sheaf of rational functions.  Hence it is a flag ideal.
Global generation of $\cO_{\cX}(D)$ relative to $\mu$ means precisely
that the evaluation map
$\mu^\ast\mu_\ast\cO_{\cX}(D)\to\cO_{\cX}(D)$ is surjective, that is,
$\mathfrak a\cO_{\cX}=\cO_{\cX}(D)$.  Applying
$\sigma(v)$ to both sides gives $\phi_{\mathfrak a}=\phi_D$, as in the
previous paragraph.

(3)
In the first direction, $\mathfrak b:=\mu_\ast\cO(D)$ contains
$\mathfrak a$.  A centered divisorial valuation on a component of
$X_{\triv}$ lifts to $\cX$ by properness.  Since
$\mathfrak a\cO_{\cX}=\cO(D)$ there, every local $f\in\mathfrak b$
has value at least that of $\mathfrak a$.  After a common
$\varpi$-shift clears the fractional denominators,
\Cref{thm:appA-valuative-integral-closure} applied in each minimal-prime
quotient makes $f$ integral over the ideal on that component.
Lift those finitely many monic integral equations to the reduced
ambient ring and multiply them.  Their evaluations belong to the
product of its minimal primes, hence vanish; their coefficients retain
the required powers of the ideal.  This proves integrality in the
ambient ring itself.  Thus
$\overline{\mathfrak b}=\overline{\mathfrak a}$.
For integral $X$ this also follows from
$\phi_{\mathfrak b}=\phi_D=\phi_{\mathfrak a}$ and
\Cref{prop:ch33-flag-faithful}.  In the second direction, $\cO_{\cX}(D)$
is the invertible pullback of $\mathfrak a$, so the universal property of
the blowup in \Cref{prop:appA-blowup-proj} yields a unique morphism
$\cX\to\operatorname{Bl}_{\mathfrak a}X_{\triv}$ over $X_{\triv}$ with
$D$ pulled back from the exceptional divisor.  Pass to the integral closure of $\cX$ in its generic fibre.  The resulting
morphism lifts uniquely through the finite relative integral closure of the
blowup, by its universal property and
\Cref{prop:ch33-integral-closure-exists}\,(2), giving the domination.
\end{proof}

\begin{remark}
\label{rem:ch33-semiample-hypothesis}
Part (2) needs relative global
generation, not merely relative semiampleness of $D$; one passes from one
to the other by replacing $D$ with a multiple, which replaces
$\phi_D$ by a positive multiple.  This is exactly why the flag-ideal
description of \Cref{chap:bj-fs-pl-psh} is by
$m^{-1}\phi_{\mathfrak a}$ rather than by $\phi_{\mathfrak a}$, and why
the resulting cone is $\PL^+(X)$ rather than a group.  By
\Cref{lem:flag-semiring-identities} the assignment is additive, so the
rescaling is recorded by the corresponding positive multiple.
\end{remark}

\begin{definition}
\label{def:ch33-rees-valuations}
For a flag ideal $\mathfrak a$, the set $\Sigma_{\mathfrak a}\subset
X^{\divis}$ of \emph{Rees valuations} of $\mathfrak a$ is the finite set
$\{v_E\}$, where $E$ runs over the irreducible components of the central
fibre of the integrally closed test configuration
$\cX_{\mathfrak a}$ of
\Cref{thm:ch33-flag-tc-equivalence}\,(1).
\end{definition}

Since the blowup of an ideal is canonically isomorphic to the blowup of any
of its powers, $\Sigma_{\mathfrak a^r}=\Sigma_{\mathfrak a}$ for
$r\geq1$.

\begin{lemma}\label{lem:ch33-rees-valuations-cover}
Every divisorial valuation on $X$ is a Rees valuation of some flag ideal;
more generally, for every finite $\Sigma\subset X^{\divis}$ there is a
flag ideal $\mathfrak a$ with $\Sigma\subset\Sigma_{\mathfrak a}$.
\end{lemma}

\begin{proof}
For each member of the finite set $\Sigma$, choose an integrally closed
configuration with a component realizing it, by
\Cref{lem:ch33-gauss-component}.  Choose a common domination of these
configurations and $X_{\triv}$, and take its relative integral closure.
Call the result $\cX$.  Every selected component has a strict transform
on $\cX$ with the same divisorial valuation: a projective birational
morphism to a scheme regular at a codimension-one point is an
isomorphism near that point.  The $\mathbf G_m$-actions and punctured
identifications preserve the normalized restrictions.

By \Cref{lem:ch33-projective-blowup-presentation}, $\cX$ is the relative
integral closure of the blowup of one flag ideal $\mathfrak a$.
Thus $\Sigma\subset\Sigma_{\mathfrak a}$.  This is the realization
argument of \cite[Lemma~2.12]{BJ22}; it makes no unsupported assertion
about unions of Rees-valuation sets under products of ideals.
\end{proof}

\begin{lemma}\label{lem:ch33-projective-blowup-presentation}
Let $\mu:\cX\to X_{\triv}$ be a test configuration.  There is a
$\mu$-ample vertical $\mathbf Q$-Cartier divisor $D$; after taking a
multiple and setting $\mathfrak a=\mu_*\cO_{\cX}(D)$, the induced map is
\[
                 \cX\longrightarrow\operatorname{Bl}_{\mathfrak a}X_{\triv}.
\]
If $\cX$ is integrally closed, it is the integral closure of this blowup
in the common generic fibre.
\end{lemma}

\begin{proof}
Write $\mu:\cX\to X_{\triv}$.  By \cite[Lemma~1.8]{BJ22}, restated: since
$\pi:\cX\to\mathbf A^1$ is projective and equivariant, there is a
$\mathbf G_m$-linearized $\pi$-ample line bundle $\cM$ on $\cX$; letting
$M$ be its restriction to $\cX_1\simeq X$ and $M_{\triv}$ the pullback
of $M$ to $X_{\triv}$, the difference
$D:=\cM-\mu^\ast M_{\triv}$ is a vertical $\mathbf Q$-Cartier divisor
which is $\mu$-ample.  After replacing $D$ by a positive multiple,
$\cO_{\cX}(D)$ is $\mu$-globally generated, so
\Cref{thm:ch33-flag-tc-equivalence}\,(2)
produces a flag ideal $\mathfrak a$ with
$\mathfrak a\cO_{\cX}=\cO_{\cX}(D)$ invertible.  The universal property
of the blowup gives a morphism $\cX\to\operatorname{Bl}_{\mathfrak a}
X_{\triv}$, which is proper and birational; since $D$ is
$\mu$-ample and pulled back from the exceptional divisor, relative
$\Proj$ reconstruction identifies $\cX$ with the integral closure of
that blowup when $\cX$ is integrally closed.  The direction of the
displayed morphism is important: $\cX$ dominates the blowup.
\end{proof}

\section{Graded norms and multiplicative filtrations}
\label{sec:ch33-norms}

We now pass from a degeneration to linear algebra.  In the section-ring
dictionary below assume $X$ integral and $L$ ample; the preceding
test-configuration and flag-ideal constructions also apply to reduced
projective $X$ with several components.  A decreasing
$\mathbf R$-filtration of $R_m$ is a family
$\mathcal F^\lambda R_m$ with
\begin{equation}
\label{eq:decreasing-filtration}
 \lambda\leq\mu\Longrightarrow
 \mathcal F^\lambda R_m\supset\mathcal F^\mu R_m,
 \quad
 \bigcup_\lambda\mathcal F^\lambda R_m=R_m,
 \quad
 \bigcap_\lambda\mathcal F^\lambda R_m=0.
\end{equation}
It is multiplicative when
\begin{equation}
\label{eq:multiplicative-filtration}
 \mathcal F^\lambda R_m\cdot\mathcal F^\mu R_{m'}
       \subset\mathcal F^{\lambda+\mu}R_{m+m'}.
\end{equation}

Equivalently, define the additive weight
\begin{equation}
\label{eq:additive-graded-norm}
 \chi_m(s):=\sup\{\lambda\mid s\in\mathcal F^\lambda R_m\},
 \qquad \chi_m(0)=+\infty .
\end{equation}
Then
\begin{equation}
\label{eq:graded-norm-inequalities}
 \chi_m(s+s')\geq\min\{\chi_m(s),\chi_m(s')\},\qquad
 \chi_{m+m'}(ss')\geq\chi_m(s)+\chi_{m'}(s').
\end{equation}
The multiplicative norm is $\|s\|_\chi=\rme^{-\chi_m(s)}$.  Since $k$ is
trivially valued, scalar multiplication by $k^\times$ does not change
$\chi_m$.

\subsection{Filtrations, norms, and jumping numbers}
\label{sec:ch33-norm-dictionary}

The passage between \eqref{eq:decreasing-filtration} and
\eqref{eq:additive-graded-norm} is a bijection once the filtration is
normalized to be left continuous, and it is worth doing carefully, because
every later argument uses one of the two languages without comment.

\begin{lemma}
\label{lem:ch33-filtration-norm-bijection}
Let $W$ be a finite-dimensional $k$-vector space.
\begin{enumerate}
\item If $(\mathcal F^\lambda W)_{\lambda\in\mathbf R}$ satisfies
 \eqref{eq:decreasing-filtration}, then
 $\chi(s):=\sup\{\lambda\mid s\in\mathcal F^\lambda W\}$ is finite for
 $s\ne0$, satisfies $\chi(as)=\chi(s)$ for $a\in k^\times$ and
 $\chi(s+s')\geq\min\{\chi(s),\chi(s')\}$, and
 $\{s\mid\chi(s)\geq\lambda\}=\bigcap_{\mu<\lambda}\mathcal F^\mu W$ is
 the left-continuous modification of $\mathcal F$.
\item Conversely such a $\chi$ determines the left-continuous filtration
 $\mathcal F^\lambda W:=\{s\mid\chi(s)\geq\lambda\}$, and the two
 constructions are mutually inverse on left-continuous filtrations.
\item For the left-continuous modification in (1), the function
 $\lambda\mapsto\dim\mathcal F^\lambda W$ is decreasing,
 left continuous, integer valued, and takes the value $\dim W$ for
 $\lambda\ll0$ and $0$ for $\lambda\gg0$.  There are exactly
 $\dim W$ jumping numbers, counted with the multiplicities
 $\dim\operatorname{gr}^\lambda_{\mathcal F}W$, where
 $\operatorname{gr}^\lambda_{\mathcal F}W
 =\mathcal F^\lambda W/\bigcup_{\mu>\lambda}\mathcal F^\mu W$.
\item There is a basis $e_1,\dots,e_N$ of $W$ adapted to
 that left-continuous filtration, meaning that $\mathcal F^\lambda W$ is spanned by those
 $e_i$ with $\chi(e_i)\geq\lambda$, for every $\lambda$.  For such a
 basis, $\chi\bigl(\sum a_ie_i\bigr)=\min\{\chi(e_i)\mid a_i\ne0\}$.
\end{enumerate}
\end{lemma}

\begin{proof}
(1)  Finiteness of $\chi(s)$ for $s\ne0$ is the union and intersection
condition in \eqref{eq:decreasing-filtration}.  Invariance under
$k^\times$ holds because each $\mathcal F^\lambda W$ is a subspace.  If
$\chi(s),\chi(s')\geq\lambda$ then $s,s'\in\mathcal F^\mu W$ for all
$\mu<\lambda$, hence so is $s+s'$, giving $\chi(s+s')\geq\lambda$.
The description of the superlevel set is immediate from the definition of a
supremum.

(2)  Clear from (1).

(3)  Monotonicity is inherited.  Left continuity holds because
$\mathcal F^\lambda W=\bigcap_{\mu<\lambda}\mathcal F^\mu W$ and a
decreasing intersection of subspaces of a finite-dimensional space
stabilizes.  The count of jumping numbers is the telescoping identity
$\dim W=\sum_\lambda\dim\operatorname{gr}^\lambda_{\mathcal F}W$, the sum being finite.

(4)  List the distinct jumping numbers as
$\lambda_1>\cdots>\lambda_q$, and put
$W_i=\mathcal F^{\lambda_i}W$, with $W_0=0$.  Choose successively a
basis of each quotient $W_i/W_{i-1}$ and lift it to $W_i$.  The union
of these lifts is a basis of $W$, and for every $i$ the vectors chosen
through stage $i$ form a basis of $W_i$.  It is therefore adapted to
every filtered subspace.  In such a basis, a nonzero linear combination
lies in $\mathcal F^\lambda W$ exactly when every basis vector occurring
in it does, which gives the displayed formula.
\end{proof}

We write
\begin{equation}
\label{eq:ch33-jumping-list}
 \lambda_{m,1}\geq\lambda_{m,2}\geq\cdots\geq\lambda_{m,N_m},
 \qquad N_m=\dim_kR_m ,
\end{equation}
for the jumping numbers of $\mathcal F$ on $R_m$, repeated with
multiplicity, as produced by
\Cref{lem:ch33-filtration-norm-bijection}\,(3).

\begin{definition}
\label{def:linearly-bounded-filtration}
A filtration is linearly bounded if there are $C_-,C_+\in\mathbf R$ such
that, for all sufficiently divisible $m$,
\begin{equation}
\label{eq:linear-bounds-filtration}
 \mathcal F^{mC_-}R_m=R_m,\qquad
 \mathcal F^{mC_+}R_m=0.
\end{equation}
\end{definition}

\begin{lemma}
\label{lem:ch33-linear-bounded-equivalents}
For a multiplicative filtration of $R(X,L)$ with $L$ ample and $X$
integral, the following are equivalent:
\begin{enumerate}
\item $\mathcal F$ is linearly bounded;
\item $\sup_{m}\,m^{-1}\max_j|\lambda_{m,j}|<\infty$;
\item $\lambda_{\max}(m):=\lambda_{m,1}$ satisfies
 $\lambda_{\max}(m)\leq Cm$ for some $C$ and all $m$, and
 $\mathcal F^{-Cm}R_m=R_m$ for the same $C$.
\end{enumerate}
Moreover, for a multiplicative filtration the limits
\begin{equation}
\label{eq:ch33-fekete-limits}
 \lambda_{\max}:=\lim_{m\to\infty}\frac{\lambda_{\max}(m)}{m}
 =\sup_m\frac{\lambda_{\max}(m)}{m}
\end{equation}
exist in $(-\infty,+\infty]$, with indices restricted to nonzero
$R_m$.  The limit is finite exactly when the filtration is linearly
bounded above.  In this full ample section ring, a uniform lower bound
already follows from multiplicativity and finite generation of
$R(X,L)$, by \cite[Lemma~1.5]{BC11}.
\end{lemma}

\begin{proof}
For any fixed retained degree, (2) and (3) bound its jumping numbers
above and below; conversely choose $C$ slightly larger than the bound
in (2), so that the strict zero condition at the upper cutoff holds.
If (1) is initially imposed only after a further Veronese change,
the upper bound extends to every retained degree: a nonzero $s\in R_m$
has $s^r\ne0$ and $\chi_{rm}(s^r)\geq r\chi_m(s)$, so a bound in degree
$rm$ bounds $\chi_m(s)$.  For the lower bound, choose finitely many
homogeneous algebra generators of the ample section ring.  Their finite
weights divided by their degrees have a common lower bound, and
multiplicativity and the ultrametric inequality extend it to every
homogeneous polynomial in those generators.  This is the proof of
\cite[Lemma~1.5]{BC11}.  Thus (1)--(3) are equivalent.

Multiplicativity gives
$\lambda_{\max}(m+m')\geq\lambda_{\max}(m)+\lambda_{\max}(m')$
whenever both section spaces are nonzero.  Reindex by the fixed
denominator $q$.  All sufficiently large reindexed spaces are nonzero,
since $L$ is ample.  For any fixed nonzero degree $p$ write a large
degree as $\ell p+r$, choosing $r$ in one fixed interval of $p$
sufficiently large degrees.  Superadditivity gives
$\lambda_{\max}(\ell p+r)\geq\ell\lambda_{\max}(p)+\lambda_{\max}(r)$.
The finite collection of remainder weights proves that the lower
limit of the normalized maxima is at least
$\lambda_{\max}(p)/p$.  Taking the supremum over $p$, and using the
opposite inequality from the definition of that supremum, proves
\eqref{eq:ch33-fekete-limits}, including its possible infinite value.
\end{proof}

Finite generation implies linear boundedness.  The converse is false: a
general filtration may have infinitely many new generators.  It is therefore
important not to identify an arbitrary graded norm with a test
configuration.  The precise statements are
\Cref{thm:ch33-fg-filtration-tc} and \Cref{ex:ch33-harmonic-not-fg}.

\subsection{From a test configuration to a filtration}
\label{sec:ch33-tc-to-filtration}

A test configuration determines such a filtration.  After clearing
denominators, a section $s\in H^0(X,mL)$ belongs to
$\mathcal F^\lambda R_m$ precisely when
$\varpi^{-\lambda}\bar s$ extends to a regular section of $m\cL$ on
$\cX$.  Equivariant multiplication proves
\eqref{eq:multiplicative-filtration}.  We state this as a definition, with
the rounding convention that makes it work for real $\lambda$.

\begin{definition}
\label{def:ch33-tc-filtration}
Let $(\cX,\cL)$ be a test configuration for $(X,L)$ and let $m$ be
sufficiently divisible.  In this definition first replace $(\cX,\cL)$ by
its integral closure in the generic fibre; this makes the convention
invariant under equivalence.  Under
\eqref{eq:tc-punctured-trivialization}, a section $s\in R_m$ determines an
equivariant rational section $\bar s$ of $m\cL$ over $\cX$, regular
away from $\cX_0$.  For $\lambda\in\mathbf R$ set
\begin{equation}
\label{eq:ch33-tc-filtration}
 \mathcal F^\lambda R_m
 :=\bigl\{s\in R_m\ \bigm|\
   \varpi^{-\lceil\lambda\rceil}\bar s\in
   H^0(\cX,m\cL)\bigr\}.
\end{equation}
\end{definition}

\begin{theorem}
\label{thm:ch33-tc-filtration-properties}
Let $(\cX,\cL)$ be a test configuration for $(X,L)$ with $X$ integral
projective and $L$ ample.
\begin{enumerate}
\item \eqref{eq:ch33-tc-filtration} defines a decreasing, left-continuous,
 $\mathbf Z$-graded filtration of each $R_m$, and the associated weight
 $\chi_m$ takes integer values.
\item The filtration is multiplicative.
\item If $\cL$ is relatively semiample it is linearly bounded, and one may
 take $C_-$ at most the smallest weight-to-degree ratio of a finite
 homogeneous generating set and $C_+$ strictly greater than its largest
 such ratio.
\item The filtration only depends on the equivalence class of
 $(\cX,\cL)$: it is unchanged by pullback along a morphism of test
 configurations and by passage to the integral closure.
\item As a $k[\varpi]$-module inside $R_m\otimes_kk[\varpi^{\pm1}]$,
\begin{equation}
\label{eq:ch33-h0-weight-decomposition}
 H^0(\cX,m\cL)
 =\bigoplus_{\lambda\in\mathbf Z}
    \mathcal F^\lambda R_m\cdot\varpi^{-\lambda},
\end{equation}
 the summand indexed by $\lambda$ being the weight-$\lambda$ subspace
 for the $\mathbf G_m$-action.
\end{enumerate}
\end{theorem}

\begin{proof}
(1)  If $\lambda\leq\mu$ then $\lceil\lambda\rceil\leq\lceil\mu\rceil$,
and $\varpi^{-\lceil\mu\rceil}\bar s
=\varpi^{-(\lceil\mu\rceil-\lceil\lambda\rceil)}
 \varpi^{-\lceil\lambda\rceil}\bar s$; multiplying a regular section by a
negative power of $\varpi$ is a stronger requirement, so
$\mathcal F^\mu R_m\subset\mathcal F^\lambda R_m$.  Each
$\mathcal F^\lambda R_m$ is a $k$-subspace because
$H^0(\cX,m\cL)$ is.  The union over $\lambda\ll0$ is all of $R_m$
because $\bar s$ is a rational section with poles only along the finitely
many components of $\cX_0$, so a sufficiently large power of $\varpi$
clears them; the intersection over $\lambda\gg0$ is zero because
$\varpi^{-N}\bar s$ has a pole along every component of $\cX_0$ for
$N$ large.  Since $\lceil\cdot\rceil$ is left continuous and integer
valued, so is $\mathcal F$, and $\chi_m$ takes integer values.

(2)  If $\varpi^{-\lambda}\bar s$ and $\varpi^{-\mu}\bar s'$ are regular
sections of $m\cL$ and $m'\cL$, their product is a regular section of
$(m+m')\cL$ and equals $\varpi^{-(\lambda+\mu)}\overline{ss'}$, because
the trivialization \eqref{eq:tc-punctured-trivialization} is compatible with
multiplication of sections.  This is
\eqref{eq:multiplicative-filtration} for integer $\lambda,\mu$, and the
general case follows from
$\lceil\lambda\rceil+\lceil\mu\rceil\geq\lceil\lambda+\mu\rceil$.

(3)  Relative semiampleness makes the section algebra
$\bigoplus_mH^0(\cX,m\cL)$ finitely generated over $k[\varpi]$ after a
Veronese change.  Choose finitely many equivariant homogeneous generators,
of degrees $d_i>0$ and weights $\lambda_i$.  Every homogeneous section
is a $k[\varpi]$-linear combination of monomials in these generators.
The ratios $\lambda_i/d_i$, together with finitely many low degrees, give
uniform upper and lower linear bounds.  This proves
\eqref{eq:linear-bounds-filtration} without assuming generation in degree
one.

(4)  By the convention in \Cref{def:ch33-tc-filtration}, both sides are
first replaced by their common integrally closed representative.  If
$\mu:\cX'\to\cX$ is then a morphism and $\cL'=\mu^\ast\cL$, then
$H^0(\cX',m\cL')=H^0(\cX,m\cL)$ because
$\mu_\ast\cO_{\cX'}=\cO_{\cX}$ whenever $\cX$ is integrally closed
(\Cref{prop:ch33-integral-closure-exists}\,(4)) and the projection formula
applies.  Hence the filtration is unchanged.  Since equivalence
of test configurations is generated by such pullbacks, (4) follows.

(5)  Write $u\in R_m\otimes_kk[\varpi^{\pm1}]$ as
$u=\sum_\lambda s_\lambda\varpi^{-\lambda}$, a finite sum with
$s_\lambda\in R_m$.  This is the weight decomposition of
\Cref{lem:ch33-weight-decomposition} for the equivariant module
$R_m\otimes k[\varpi^{\pm1}]$, the weight of $s\varpi^{-\lambda}$ being
declared to be $\lambda$.  By \Cref{lem:ch33-weight-decomposition} again,
$u$ is a section of $m\cL$ over $\cX$ if and only if every weight
component is, and $s_\lambda\varpi^{-\lambda}\in H^0(\cX,m\cL)$ means
exactly $s_\lambda\in\mathcal F^\lambda R_m$.
\end{proof}

\begin{remark}
\label{rem:ch33-weight-sign}
Assertion (5) is the definition of the weight labeling used throughout the
book: the weight of $s\varpi^{-\lambda}$ is $\lambda$, so that the
weight-$\lambda$ subspace of $H^0(\cX,m\cL)$ is
$\mathcal F^\lambda R_m\varpi^{-\lambda}$.  With this labeling, larger
weight means deeper in the filtration, and the twist rule of
\Cref{prop:ch33-transformations} holds with a plus sign.  A source that lets
a one-parameter subgroup act on \emph{points} rather than on homogeneous
coordinates obtains the opposite sign; the reflection is detected
immediately by the position of the support of the Duistermaat--Heckman
measure, and this is one of the unit tests carried out in
\Cref{ex:p1-product-dh}.
\end{remark}

\begin{proposition}
\label{prop:ch33-central-fibre-weights}
Let $(\cX,\cL)$ be a test configuration for $(X,L)$ with $\cL$
relatively ample and $\cX$ integrally closed, and let $\cL_0$ be the
restriction of $\cL$ to $\cX_0$.  Then for all sufficiently divisible
$m\gg0$:
\begin{enumerate}
\item $H^0(\cX,m\cL)$ is a free $k[\varpi]$-module of rank
 $N_m=h^0(X,mL)$;
\item the restriction map
 $H^0(\cX,m\cL)\to H^0(\cX_0,m\cL_0)$ is surjective with kernel
 $\varpi H^0(\cX,m\cL)$, so that
\begin{equation}
\label{eq:ch33-central-fibre-graded}
 H^0(\cX_0,m\cL_0)
 \simeq
 \bigoplus_{\lambda\in\mathbf Z}\operatorname{gr}^\lambda_{\mathcal F}R_m ,
\end{equation}
 the summand indexed by $\lambda$ being the weight-$\lambda$ subspace;
\item consequently the weights of the $\mathbf G_m$-representation
 $H^0(\cX_0,m\cL_0)$, listed with multiplicity, are exactly the jumping
 numbers $\lambda_{m,1}\geq\cdots\geq\lambda_{m,N_m}$ of
 \eqref{eq:ch33-jumping-list}, and the total weight is
\begin{equation}
\label{eq:ch33-total-weight}
 w_m:=\sum_{\lambda}\lambda\dim\operatorname{gr}^\lambda_{\mathcal F}R_m
     =\sum_{j=1}^{N_m}\lambda_{m,j}.
\end{equation}
\end{enumerate}
\end{proposition}

\begin{proof}
(1)  Properness makes $\pi_\ast(m\cL)$ a coherent
$k[\varpi]$-module.  It is torsion free, being a submodule of
$R_m\otimes k[\varpi^{\pm1}]$ by
\eqref{eq:ch33-h0-weight-decomposition}, hence free, $k[\varpi]$ being a
principal ideal domain.  Its rank is the dimension of its generic fibre,
which is $h^0(X_{k(\varpi)},mL_{k(\varpi)})=N_m$.

(2)  Serre vanishing for the $\pi$-ample $\cL$ gives
$R^1\pi_\ast(m\cL)=0$ for $m\gg0$.  It makes
$\pi_\ast(m\cL)\otimes_{k[\varpi]}k[\varpi]/(\varpi)
\to H^0(\cX_0,m\cL_0)$ an isomorphism, by cohomology and base change for
the flat morphism $\pi$.  Reducing
\eqref{eq:ch33-h0-weight-decomposition} modulo $\varpi$ and using
$\varpi\cdot\mathcal F^\lambda R_m\varpi^{-\lambda}
=\mathcal F^\lambda R_m\varpi^{-(\lambda-1)}\subset
\mathcal F^{\lambda-1}R_m\varpi^{-(\lambda-1)}$ identifies the quotient
with $\bigoplus_\lambda(\mathcal F^\lambda R_m/\mathcal F^{\lambda+1}R_m)$,
which is \eqref{eq:ch33-central-fibre-graded}.

(3)  Immediate from (2) and
\Cref{lem:ch33-filtration-norm-bijection}\,(3), the sum of the
multiplicities being $N_m$ by (1).
\end{proof}

\begin{remark}
\label{rem:ch33-ampleness-needed}
Only assertion (2) uses relative ampleness, and it uses it only through the
vanishing of $R^1\pi_\ast$.  For a merely semiample $\cL$ the
filtration of \Cref{def:ch33-tc-filtration} is still defined and still
multiplicative, and \eqref{eq:ch33-h0-weight-decomposition} still holds; what
may fail is the identification of $\operatorname{gr}_{\mathcal F}R_m$ with sections on
the central fibre.  The Duistermaat--Heckman measure of
\Cref{sec:ch33-dh} is defined from the filtration rather than from the
central fibre, so only the interpretation
\eqref{eq:ch33-central-fibre-graded} is lost.
\end{remark}

\subsection{From a filtration to a test configuration}
\label{sec:ch33-filtration-to-tc}

Conversely, a finitely generated rational filtration has Rees algebra
\begin{equation}
\label{eq:filtration-rees}
 \mathcal R(\mathcal F)
   :=\bigoplus_{m\geq0}\ \bigoplus_{\lambda\in\mathbf Z}
       \mathcal F^\lambda R_m\,\varpi^{-\lambda},
\end{equation}
and its relative $\Proj$ gives a test configuration after taking integral
closure.  Changing all weights by $cm$ twists the linearization by the
character of weight $c$; analytically it adds the constant $c$.

\begin{definition}
\label{def:ch33-fg-filtration}
A multiplicative filtration $\mathcal F$ of $R(X,L)$ is
\emph{rational}, in the integral-weight convention of this Rees subsection,
if all its jumping numbers are integers after restriction to a suitable
Veronese section ring.  It is \emph{finitely generated} if the bigraded algebra
\eqref{eq:filtration-rees} is a finitely generated $k$-algebra.
\end{definition}

We normalize degree zero by
$\mathcal F^\lambda R_0=k$ for $\lambda\leq0$ and $0$ for $\lambda>0$;
this is the convention used in \eqref{eq:filtration-rees}.
Then \eqref{eq:filtration-rees} contains
$k[\varpi]$ in bidegree $(0,\ast)$, and its localization at
$\varpi$ is $R(X,L)\otimes_kk[\varpi^{\pm1}]$; the next theorem shows
that finite generation upgrades these two elementary
observations to a geometric statement.

\begin{theorem}
\label{thm:ch33-fg-filtration-tc}
Let $X$ be integral projective and $L$ ample.
\begin{enumerate}
\item If $(\cX,\cL)$ is a test configuration for $(X,L)$ with $\cL$
 relatively semiample and rational, then the filtration
 \eqref{eq:ch33-tc-filtration} is multiplicative, linearly bounded,
 rational, and finitely generated.
\item Conversely, let $\mathcal F$ be a multiplicative, linearly bounded,
 rational, finitely generated filtration.  Then for a suitable Veronese
 index $r$, the relative $\Proj$ over $\mathbf A^1$ of the
 $\varpi$-graded algebra \eqref{eq:filtration-rees} in Veronese degrees
 divisible by $r$, followed by integral closure in its generic fibre, is
 an integrally closed test configuration $(\cX_{\mathcal F},\cL_{\mathcal
 F})$ for $(X,rL)$ with $\cL_{\mathcal F}$ relatively semiample.  Its
 filtration is the Rees/irrelevant saturation of $\mathcal F$, followed
 by integral closure; these operations may enlarge the original filtered
 pieces.
\item The two constructions are mutually inverse: starting from a
 relatively semiample $(\cX,\cL)$ with $\cX$ integrally closed and
 applying (1) then (2) returns a test configuration equivalent to
 $(\cX,\cL)$.  Starting from $\mathcal F$ and applying (2) then (1)
 returns $\mathcal F$ precisely when its Veronese Rees algebra is already
 irrelevant-saturated and integrally closed.
\end{enumerate}
Thus, with the identifications above, integrally closed semiample test
configurations for $(X,L)$ modulo equivalence and Veronese change
correspond bijectively to multiplicative, linearly bounded, rational,
finitely generated filtrations of $R(X,L)$ with a fixed linearization,
modulo passage to a common saturated integral closure of their Rees
algebras.
\end{theorem}

\begin{proof}
(1) Multiplicativity, linear boundedness, and rationality are \Cref{thm:ch33-tc-filtration-properties}\,(1)--(3). For finite generation, take the ample model of the relatively semiample configuration. By \cite[Proposition~2.17]{BHJ17}, the polarization is pulled back from an ample one and the contraction $\mu$ satisfies $\mu_*\mathcal O=\mathcal O$. The projection formula identifies their section algebras; the ample model has a finitely generated section algebra after clearing denominators. By \eqref{eq:ch33-h0-weight-decomposition}, this is the Rees algebra. Relative semiampleness gives this contraction, rather than an embedding of the original total space.

(2)  Let $\mathcal A=\mathcal R(\mathcal F)$.  Finite generation lets us
choose a Veronese index $r$ such that the $r$-th Veronese subalgebra
$\mathcal A^{(r)}=\bigoplus_{m\geq0}\mathcal A_{rm}$ is generated in
degree one over $\mathcal A_0=k[\varpi]$; this is the standard fact that
a finitely generated graded algebra becomes generated in degree one after
passing to a Veronese subalgebra.  Set
$\cX_{\mathcal F}:=\Proj_{\mathbf A^1}\mathcal A^{(r)}$ with the
tautological $\cO(1)=\cL_{\mathcal F}$.  Then:
\emph{projectivity} holds because $\mathcal A^{(r)}$ is generated in
degree one by a finite-dimensional $k[\varpi]$-module;
\emph{flatness} holds because each graded piece
$\mathcal A_{rm}=\bigoplus_\lambda\mathcal F^\lambda R_{rm}
\varpi^{-\lambda}$ is a torsion-free $k[\varpi]$-module, being a
submodule of $R_{rm}\otimes k[\varpi^{\pm1}]$;
the \emph{$\mathbf G_m$-action} is the second grading, and it lifts the
standard action on $\mathbf A^1=\Spec k[\varpi]$ because
$\varpi$ has weight $-1$ in \eqref{eq:filtration-rees};
the \emph{trivialization} \eqref{eq:tc-punctured-trivialization} comes from
$\mathcal A^{(r)}[\varpi^{-1}]=R^{(r)}(X,L)\otimes k[\varpi^{\pm1}]$,
whose relative $\Proj$ is $X\times\mathbf G_m$;
and $\cL_{\mathcal F}$ is relatively semiample by construction, restricting
to $rL$ over $1$.  Taking integral closure in the generic fibre changes
none of this, by
\Cref{prop:ch33-integral-closure-exists}\,(1).  Finally, by
\eqref{eq:ch33-h0-weight-decomposition} applied to
$\cX_{\mathcal F}$, its sections in degree $m$ are the irrelevant
saturation of $\mathcal A_{rm}$.  Passing to the integral closure can
enlarge them further.  Thus the returned filtration is exactly the
saturated integral-closure filtration stated in (2), not necessarily the
original $\mathcal F$.

(3) For the first composite, (1) identifies the Rees algebra with the section algebra. Its relative $\Proj$ gives the ample model, which is equivalent to the original semiample configuration and is integrally closed; it need not be the original total space. This is \cite[Proposition~2.17]{BHJ17} and \cite[Proposition~2.30]{BJ22}. For the second composite, (2) gives the original filtration exactly under the stated saturation and integral-closure hypotheses. Uniqueness of the ample integrally closed representative, in the cited proposition of \cite{BJ22}, gives the final bijection.
\end{proof}

\begin{warning}{Veronese change is not optional}
\label{warn:ch33-veronese}
The bijection of \Cref{thm:ch33-fg-filtration-tc} is stated modulo Veronese
change, and it must be.  A finitely generated bigraded algebra need not be
generated in degree one.  Its relative $\Proj$ already exists, while
a suitable Veronese ensures that the tautological polarization is an
honest line bundle generated in degree one.  Different choices give canonically
isomorphic test configurations, but the polarization changes by the
corresponding multiple.  In \Cref{chap:bj-stability} this is what allows
norms to be identified only after passing to a common Veronese index.
\end{warning}

\begin{proposition}
\label{prop:ch33-transformations}
Let $\mathcal F$ be a multiplicative linearly bounded filtration with
weights $\chi_m$, let $c\in\mathbf R$ and $d\in\mathbf Z_{>0}$.
\begin{enumerate}
\item \emph{Twist.} $\chi_m^{(c)}:=\chi_m+mc$ is again multiplicative and linearly bounded. If $\mathcal F$ comes from $(\cX,\cL)$ and $c$ is rational, it is induced by the character twist after clearing polarization denominators, represented after compactification by $\overline\cL+c\,\overline\pi^*\{0\}$. An arbitrary real $c$ defines the shifted filtration but need not define an algebraic character.
\item \emph{Base change.} $\chi_m^{[d]}:=d\chi_m$ is again multiplicative and linearly bounded. For a test-configuration filtration it gives the weights on the raw base change along $\varpi\mapsto\varpi^d$. Subsequent integral closure in the generic fibre can increase these weights. The resulting integral-closure filtration of \Cref{def:ch33-tc-filtration} therefore need not equal $\chi^{[d]}$, although their Fubini--Study functions agree.  When $X$ is normal, this operation is the usual normalized base change.
\item \emph{Veronese.}  The restriction of $\mathcal F$ to
 $R^{(r)}=\bigoplus_mR_{rm}=R(X,rL)$, with degree-$m$ weights
 $\chi_{rm}$, is again multiplicative and linearly bounded.  When
 $\mathcal F$ comes from $(\cX,\cL)$, it is the filtration of
 $(\cX,r\cL)$.  For an arbitrary filtration this restriction does not
 recover the filtered pieces in the omitted degrees.
\item \emph{Scaling of the analytic side.}  On the associated function of
 \Cref{sec:ch33-base-ideals} these three operations act by
 $\phi\mapsto\phi+c$, $\phi\mapsto d\cdot\phi$ in the sense of the
 scaling action \eqref{eq:valuation-scaling}, and $\phi\mapsto r\phi$
 with the class $L$ replaced by $rL$.  Equivalently, if retained
 degrees are still normalized by their actual degree $rm$, Veronese
 restriction leaves $\phi$ unchanged.
\end{enumerate}
\end{proposition}

\begin{proof}
(1) Multiplicativity follows by adding $mc$ in degree $m$. For a rational character twist, work in degrees with $mc\in\mathbb Z$. If $\bar s$ is the invariant extension before twisting, the new invariant extension is $\varpi^{mc}\bar s$. Thus $\varpi^{-\lambda}\varpi^{mc}\bar s$ extends exactly when the old extension with exponent $\lambda-mc$ does. This gives $\chi_m+mc$ with the asserted sign. The compactified description is \Cref{thm:ch33-intersection-invariance}\,(3).

(2) Multiplicativity follows by multiplying weights by $d$. Flat base change pulls an adapted section-lattice basis $\varpi^{-\lambda_j}\bar s_j$ back to one with exponents $-d\lambda_j$. Hence the raw family has weights $d\chi_m$. Taking integral closure can enlarge the section lattice, so \Cref{thm:ch33-tc-filtration-properties}\,(4) cannot identify these raw weights with the integral-closure weights: that result first takes integral closure on both sides. The graded integral-dependence argument in \Cref{lem:rf-integral-closure} proves equality of the Fubini--Study functions. \Cref{ex:rf-normalization-changes-weights} exhibits strict enlargement of weights under double base change followed by integral closure.

(3)  Restriction to the degrees divisible by $r$ preserves
multiplicativity and the linear bounds, now with constants $rC_-$ and
$rC_+$ when the new degree is $m$.  The section lattice for
$(\cX,r\cL)$ in degree $m$ is the original lattice in degree $rm$,
so its weight function is $\chi_{rm}$.  No relation in these retained
degrees specifies a filtered subspace in a degree not divisible by $r$,
which proves the final assertion.

(4)  Substituting into the formula \eqref{eq:norm-fs-function} for the
associated function gives $\phi+c$ in case (1); in case (2) it gives
$\phi_{\chi^{[d]}}(v)=\sup(d\chi_m(s)-v(s))/m
=d\sup(\chi_m(s)-d^{-1}v(s))/m=d\,\phi_\chi(d^{-1}v)$, which is the scaling
action; in case (3) the class is multiplied by $r$ and so is the
function.  For case (3), multiplicativity gives
$\chi_{rm}(s^r)\geq r\chi_m(s)$, while valuations satisfy
$v(s^r)=rv(s)$.  Thus restricting the supremum to divisible degrees does
not change it when the denominator remains the actual degree $rm$; after
reindexing $rm$ as degree $m$ for $rL$, the value is multiplied by
$r$.
\end{proof}

\section{Base ideals and the associated analytic function}
\label{sec:ch33-base-ideals}

Given a filtered linear series, the image of
\[
 \mathcal F^\lambda R_m\otimes\cO_X(-mL)\longrightarrow\cO_X
\]
is its base ideal $\mathfrak b_{m,\lambda}$.  Multiplicativity gives
\begin{equation}
\label{eq:filtered-base-ideals}
 \mathfrak b_{m,\lambda}\mathfrak b_{m',\mu}
       \subset\mathfrak b_{m+m',\lambda+\mu}.
\end{equation}
For a test configuration, these ideals recover its flag ideal in each
degree.  The associated function can be read in either language:
\begin{equation}
\label{eq:norm-fs-function}
 \phi_\chi(v)
   =\sup_{\substack{m\geq1\\0\ne s\in R_m}}
        \frac{\chi_m(s)-v(s)}{m}.
\end{equation}
Here $v(s)$ is defined by choosing a local frame of $mL$ at $c(v)$.
Two frames differ by a unit $u$, and domination by the valuation ring gives
$v(u)=v(u^{-1})=0$; hence the value is independent of the frame.  If a
section vanishes identically on the support $s(v)$, its value is
$+\infty$.
For a finitely generated filtration the supremum reduces, after a Veronese
subalgebra, to a finite maximum; this is the Fubini--Study function of the
test configuration.  For the full ample section ring with a lower linear bound, the supremum is finite even at proper-support semivaluations. Indeed, a globally generated degree has a section of value zero at the center, with weight at least $mC_-$. The upper linear bound gives the opposite bound $\phi_\chi\leq C_+$. Individual terms can be $-\infty$, but their supremum lies in $[C_-,C_+]$; no cancellation of infinities occurs.

\subsection{Base ideals compute the flag ideal}
\label{sec:ch33-base-ideal-flag}

\begin{proposition}\label{prop:ch33-base-ideal-flag}
Let $\mathcal F$ be a multiplicative, linearly bounded filtration of
$R(X,L)$, with $mL$ globally generated for $m$ sufficiently
divisible, and assume that its jumping numbers in those degrees are
integers.  This includes the integral-weight rational filtrations of
\Cref{def:ch33-fg-filtration} after their Veronese restriction.
Rational numerical values alone do not imply this integral presentation.
\begin{enumerate}
\item The base ideals satisfy \eqref{eq:filtered-base-ideals}, they are
 decreasing in $\lambda$, and $\mathfrak b_{m,\lambda}=\cO_X$ for
 $\lambda\leq mC_-$ whenever
 $\mathcal F^{mC_-}R_m=R_m$.
\item In each degree $m$, the vertical fractional ideal
\begin{equation}
\label{eq:ch33-degree-m-flag}
 \mathfrak a_{(m)}
 :=\sum_{\lambda\in\mathbf Z}\mathfrak b_{m,\lambda}\varpi^{-\lambda}
\end{equation}
 is a flag ideal on $X_{\triv}$, and
\begin{equation}
\label{eq:ch33-flag-vs-norm}
 \frac1m\phi_{\mathfrak a_{(m)}}(v)
 =\frac1m\max_{\lambda}\bigl\{\log|\mathfrak b_{m,\lambda}|(v)+\lambda\bigr\}
 =\sup_{0\ne s\in R_m}\frac{\chi_m(s)-v(s)}{m}
\end{equation}
 for every $v\in X^{\An}$.
\item Consequently
\begin{equation}
\label{eq:ch33-phi-sup-of-flags}
 \phi_\chi=\sup_{m}\ \frac1m\,\phi_{\mathfrak a_{(m)}},
\end{equation}
 an increasing supremum along the sufficiently divisible integers.
\item If $\mathcal F$ is the filtration of a relatively semiample test
 configuration, choose a Veronese index $r$ for which its saturated Rees
 algebra is generated in degree one, and put
 $\mathfrak a=\mathfrak a_{(r)}$.  Then for every $\ell\geq1$,
 \begin{equation}
 \label{eq:ch33-flag-power-propagation}
  \overline{\mathfrak a_{(r\ell)}}
  =\overline{\mathfrak a^{\ell}} .
 \end{equation}
 Thus the flag ideal is degree-dependent before a generating Veronese is
 fixed; the exponent is $\ell$, not the original degree $r\ell$.
\end{enumerate}
\end{proposition}

\begin{proof}
(1)  Base ideals are decreasing because the filtration is.  The product
inclusion follows from \eqref{eq:multiplicative-filtration}: the image of
the product of the two evaluation maps is contained in the image of the
evaluation map in degree $m+m'$ and weight $\lambda+\mu$.  If
$\mathcal F^{mC_-}R_m=R_m$ then $\mathfrak b_{m,mC_-}$ is the base ideal
of the complete linear system $|mL|$, which is $\cO_X$ by global
generation.

(2)  Each $\mathfrak b_{m,\lambda}$ is a coherent ideal, equal to
$\cO_X$ for $\lambda\ll0$ and to $0$ for $\lambda\gg0$ by linear
boundedness, and they decrease in $\lambda$; so
\eqref{eq:ch33-degree-m-flag} is a flag ideal by
\Cref{lem:ch33-flag-weight-decomposition}.  The first equality in
\eqref{eq:ch33-flag-vs-norm} is \eqref{eq:flag-ideal-function} in the
indexing of \eqref{eq:ch33-flag-weight}.  For the second, fix
$v\in X^{\An}$.  By definition
$\log|\mathfrak b_{m,\lambda}|(v)=-v(\mathfrak b_{m,\lambda})
=\max\{-v(s)\mid s\in\mathcal F^\lambda R_m\}$, because a local generating
set of $\mathfrak b_{m,\lambda}$ at $c(v)$ is obtained from a basis of
$\mathcal F^\lambda R_m$ written in a local frame, and the value of an
ideal is the minimum over generators
\eqref{eq:ideal-evaluation-triv}.  Hence
\[
 \max_\lambda\{\log|\mathfrak b_{m,\lambda}|(v)+\lambda\}
 =\max_\lambda\max_{s\in\mathcal F^\lambda R_m\setminus0}
          \{\lambda-v(s)\}
 =\max_{0\ne s\in R_m}\{\chi_m(s)-v(s)\},
\]
the last step because the jumping numbers are integral, so for fixed
$s\ne0$ the largest admissible integer $\lambda$ is $\chi_m(s)$.

(3)  Combine (2) with \eqref{eq:norm-fs-function}.  The supremum is
increasing along divisibility because
$\phi_{\mathfrak a_{(m)}}+\phi_{\mathfrak a_{(m')}}
\leq\phi_{\mathfrak a_{(m+m')}}$, by \eqref{eq:filtered-base-ideals} and
\Cref{lem:flag-semiring-identities}; in particular
$m^{-1}\phi_{\mathfrak a_{(m)}}\leq
(rm)^{-1}\phi_{\mathfrak a_{(rm)}}$ for $r\geq1$.

(4)  The degree-$m$ part of the Rees algebra is the section lattice
encoded by $\mathfrak a_{(m)}$; this is precisely the degree-specific
flag ideal of \cite[Definition~2.19 and Proposition~2.21]{BHJ17}.  After
passing to the chosen Veronese, generation in degree one says that the
degree-$\ell$ lattice is generated by products of $\ell$ elements of
the degree-one lattice.  Hence its flag ideal equals
$\mathfrak a^\ell$ before saturation.  Passing to the saturated
integral-closure representative replaces both sides by their integral
closures and gives \eqref{eq:ch33-flag-power-propagation}.
\end{proof}

\begin{corollary}\label{cor:ch33-fg-finite-max}
If $\mathcal F$ is rational and finitely generated, then after passing to
a Veronese index $r$ as in \Cref{thm:ch33-fg-filtration-tc}\,(2), the
supremum \eqref{eq:norm-fs-function} is attained at $m=r$, so that
\begin{equation}
\label{eq:ch33-finite-max}
 \phi_\chi(v)
 =\frac1r\max_{1\leq i\leq N}\{\chi_r(s_i)-v(s_i)\}
\end{equation}
for any basis $s_1,\dots,s_N$ of $R_r$ adapted to $\mathcal F$.  In
particular $\phi_\chi$ is a rational Fubini--Study function for $L$ in
the sense of \Cref{def:gf-fs-functions}, and it is continuous and
real-valued as soon as the sections $s_i$ have no common zero.
\end{corollary}

\begin{proof}
By \Cref{thm:ch33-fg-filtration-tc}\,(2), $\mathcal R(\mathcal F)^{(r)}$
is generated in degree one over $k[\varpi]$.  Let
$s_i\varpi^{-\chi_r(s_i)}$ be the corresponding generators, $s_i$
running over an adapted basis of $R_r$ as in
\Cref{lem:ch33-filtration-norm-bijection}\,(4).  Any element of
$\mathcal F^\lambda R_{rm}$ is a $k[\varpi]$-linear combination of
products of $m$ generators, so it is a sum of terms
$a\,s_{i_1}\cdots s_{i_m}$ with
$\sum_j\chi_r(s_{i_j})\geq\lambda$.  Evaluating a nonzero such element
$s$ at $v$ and using the ultrametric inequality
$v(\sum)\geq\min v$ gives
\[
 \chi_{rm}(s)-v(s)
 \leq\max\Bigl\{\sum_j\bigl(\chi_r(s_{i_j})-v(s_{i_j})\bigr)\Bigr\}
 \leq m\max_i\{\chi_r(s_i)-v(s_i)\},
\]
so the quotient by $rm$ never exceeds the right-hand side of
\eqref{eq:ch33-finite-max}.  The reverse inequality is the case $m=1$.
The identification with a Fubini--Study function is
\eqref{eq:gf-fs-presentation} with $\lambda_i=\chi_r(s_i)$.
\end{proof}

\begin{proposition}\label{prop:tc-valuative-invariance}
For a rational finitely generated filtration, integral closure of its
Rees algebra inside the localized generic section algebra does not change
$\phi_\chi$.  Two relatively semiample integrally closed test configurations
with fixed ample generic polarization determine the same function
$\phi_\chi$ if and only if they are equivalent.
\end{proposition}

\begin{proof}
The first assertion is the graded integral-dependence argument in
\Cref{lem:rf-integral-closure}.  A homogeneous integral equation compares
the weight and valuation of a sufficiently high power of a section with
those of sections in the original graded algebra, proving equality of
the resulting suprema.  It does not assert that graded integral closure
acts independently in each degree.

For a relatively semiample configuration, global generation in a
sufficiently divisible degree identifies $\phi_\chi$ with the function
of its vertical difference divisor on a common domination of the
product; this is \cite[Theorem~2.31(ii), proof]{BJ22}.
Choose an integrally closed common domination of the two configurations
and the product.  The difference of the pulled-back polarizations is
a vertical $\mathbf Q$-Cartier divisor $D$.  Equality of their functions
gives $\phi_D=0$.  Apply \Cref{lem:ch33-effectivity} to $D$ and $-D$:
both are effective, so $D=0$ and the configurations are equivalent.
The converse follows by pullback invariance.

For an arbitrary polarization the vertical-divisor function still
determines its equivalence class, by \cite[Theorem~2.31(i)]{BJ22};
the section supremum $\phi_\chi$ need not equal that function without
semiampleness.  This is why the second assertion has the stated scope.
\end{proof}

\subsection{A finite-generation criterion}
\label{sec:ch33-fg-criterion}

Deciding whether a given filtration is finitely generated is in general
hard; \Cref{chap:bj-stability} records the deep finite-generation theorems
of \cite{LXZ22} for filtrations computing stability thresholds.  For the
diagonalizable filtrations occurring in our examples, however, the following
semigroup criterion makes the obstruction explicit.

\begin{proposition}
\label{prop:ch33-monomial-fg-criterion}
Let $S\subset\mathbf Z^N$ be a subsemigroup containing $0$, and let
$k[S]\subset k[T_1^{\pm1},\dots,T_N^{\pm1}]$ be the corresponding
semigroup algebra.  If $k[S]$ is a finitely generated $k$-algebra, then
$S$ is a finitely generated semigroup and the convex cone
$\operatorname{Cone}(S)\subset\mathbf R^N$ generated by $S$ is
rational polyhedral.

Now let $C\subset\mathbf R^{N-1}$ be a rational polyhedral cone, let
$\Psi:C\to\mathbf R$ be continuous, concave, and positively homogeneous
of degree one, and let
\begin{equation}
\label{eq:ch33-hypograph-semigroup}
 S_\Psi:=\{(a,\lambda)\in(C\cap\mathbf Z^{N-1})\times\mathbf Z
            \mid\lambda\leq\Psi(a)\}.
\end{equation}
Then $S_\Psi$ is a subsemigroup of $\mathbf Z^N$, and
$k[S_\Psi]$ is a finitely generated $k$-algebra if and only if the
hypograph
\begin{equation}
\label{eq:ch33-hypograph}
 \operatorname{hypo}(\Psi)
 :=\{(x,\tau)\in C\times\mathbf R\mid\tau\leq\Psi(x)\}
\end{equation}
is a rational polyhedral cone, that is, if and only if $\Psi$ is the
minimum of finitely many linear forms with rational coefficients.
\end{proposition}

\begin{proof}
For the first paragraph, let $f_1,\dots,f_q$ be algebra generators of
$k[S]$, and let $S_0\subset S$ be the finite union of their monomial
supports.  For any $s\in S$, write the monomial $T^s$ as a polynomial in
the $f_i$ and expand.  Since the coefficient of $T^s$ in that expansion is
nonzero, at least one contributing product of monomials from the supports
of the $f_i$ has exponent $s$.  Thus $s$ is a sum of elements of $S_0$.
Hence $S$ is generated as a semigroup by the finite set $S_0$, and
$\operatorname{Cone}(S)=\sum_{s\in S_0}\mathbf R_{\geq0}s$ is rational
polyhedral, in particular closed.

For the second paragraph, first note that $S_\Psi$ is a semigroup:
concavity and homogeneity of $\Psi$ give superadditivity,
$\Psi(a+a')\geq\Psi(a)+\Psi(a')$, so
$\lambda+\lambda'\leq\Psi(a+a')$ whenever
$\lambda\leq\Psi(a)$ and $\lambda'\leq\Psi(a')$.

Assume $k[S_\Psi]$ is finitely generated.  By the first paragraph,
$\operatorname{Cone}(S_\Psi)$ is a rational polyhedral cone, hence
closed.  We claim it equals $\operatorname{hypo}(\Psi)$.  The inclusion
$\subset$ holds because $\operatorname{hypo}(\Psi)$ is a convex cone
containing $S_\Psi$, by concavity and homogeneity.  For $\supset$, let
$(x,\tau)\in\operatorname{hypo}(\Psi)$ be rational.  Choose
$M\in\mathbf Z_{>0}$ with $Mx\in\mathbf Z^{N-1}$ and
$M\tau\in\mathbf Z$; homogeneity gives
$M\tau\leq M\Psi(x)=\Psi(Mx)$, so $(Mx,M\tau)\in S_\Psi$ and therefore
$(x,\tau)\in\operatorname{Cone}(S_\Psi)$.  Rational points are dense in $\operatorname{hypo}(\Psi)$: first
approximate $x\in C$ by rational points of the rational cone $C$,
and then choose rational heights strictly below their $\Psi$-values,
approaching the desired height.  Continuity of $\Psi$ justifies this
approximation even at the boundary of $C$.  Density and closedness give
$\operatorname{hypo}(\Psi)\subset\operatorname{Cone}(S_\Psi)$.  Thus
$\operatorname{hypo}(\Psi)$ is rational polyhedral; writing it as a finite
intersection of rational half-spaces and using that it is a hypograph over
$C$ presents $\Psi$ as a minimum of finitely many rational linear
forms.

Conversely, if $\operatorname{hypo}(\Psi)$ is a rational polyhedral cone,
then $S_\Psi=\operatorname{hypo}(\Psi)\cap\mathbf Z^N$ is a finitely
generated semigroup by Gordan's lemma, so $k[S_\Psi]$ is a finitely
generated $k$-algebra.
\end{proof}

\begin{warning}{Irrationality is not the only obstruction}
\label{warn:ch33-not-only-irrationality}
It is tempting to believe that a filtration fails to be finitely generated
only because its jumping numbers are irrational.  This is false.  In
\Cref{ex:ch33-harmonic-not-fg} every normalized jumping number is rational,
every weight is a rational number, and yet the integer-level Rees
algebra is not finitely
generated, because the relevant hypograph is not polyhedral.  The
obstruction is polyhedrality, not rationality.
\end{warning}

\section{Duistermaat--Heckman measures}
\label{sec:ch33-dh}

Assume $X$ is integral of dimension $n\geq1$, $L$ ample, and
$\mathcal F$ linearly bounded and multiplicative.  For the omitted
zero-dimensional case $X=\Spec k$, every retained section space has
dimension one; Fekete's limit gives a single Dirac measure.  The
positive-dimensional volume and $1/n$ formulas below concern $n\geq1$.  Let
\begin{equation}
\label{eq:jumping-numbers}
 \lambda_{m,1}\geq\cdots\geq\lambda_{m,N_m},
 \qquad N_m=h^0(X,mL),
\end{equation}
be the jumping numbers of $\mathcal F$, repeated with multiplicity.  Define
\begin{equation}
\label{eq:empirical-weight-measure}
 \nu_m:=\frac1{N_m}\sum_{j=1}^{N_m}
                    \delta_{\lambda_{m,j}/m}.
\end{equation}

\refstepcounter{theorem}\label{thm:dh-limit-filtration}
\begin{quote}
\textbf{Foundational input \thetheorem\ (Boucksom--Chen convergence of
weight distributions; source boundary).}
The measures $\nu_m$ converge weakly, along the sufficiently divisible
integers, to a compactly supported probability measure
$\DHm(\mathcal F)$.  For a test configuration, this is the asymptotic
weight distribution of the $\mathbf G_m$-representation on the central
fibre.  It depends only on the corresponding graded norm.
This is the normalized measure statement of \cite[Theorem~A]{BC11}, with
the test-configuration identification supplied by
\Cref{prop:ch33-central-fibre-weights}; the convergence theorem itself is
imported and no proof is claimed here.
\end{quote}

Before recording its normalization and consequences we fix
the normalization, because three different normalizations of the same object
occur in the literature and in this book.

\begin{warning}{Three normalizations}
\label{warn:ch33-dh-normalizations}
The measure $\DHm(\mathcal F)$ of \Cref{thm:dh-limit-filtration} is a
\emph{probability} measure.  It coincides with the Duistermaat--Heckman
measure of \cite{BHJ17} and with $\DHm^{\mathrm{prob}}$ of
\Cref{thm:rf-dh-quantization}.  The \emph{unnormalized} measure of
\Cref{sec:appJ-rays}, of total mass $V=(L^n)$, is
\begin{equation}
\label{eq:ch33-dh-unnormalized}
 \DHm^{\mathrm u}(\mathcal F):=V\cdot\DHm(\mathcal F),
\end{equation}
which is the convention of \eqref{eq:cma-dh-normalization}.  Finally, the
Okounkov-body normalization of \Cref{sec:dh-preview} has total mass
$\vol(\Delta)=V/n!$.  The three differ by the positive constants $V$
and $n!$, and we always display which one is in force.  A statement about
moments is meaningless until this has been done.
\end{warning}

\subsection{Existence of the limit}
\label{sec:ch33-dh-existence}

Fix once and for all a valuation
$\nu:k(X)^\times\to\mathbf Z^n$, of rank $n$ with one-dimensional
leaves, for the lexicographic order; such valuations exist on any integral
projective variety of dimension $n$, for instance from a flag of
subvarieties through a regular closed point of a birational model
\cite{Oko03,LM09,KK12}.  For a graded linear series
$W_\bullet\subset R(X,L)$, let
$\Gamma(W_\bullet)\subset\mathbf N\times\mathbf Z^n$ be the associated
graded semigroup and $\Delta(W_\bullet)$ its Okounkov body, the compact
convex set obtained by intersecting the closed cone generated by
$\Gamma(W_\bullet)$ with the slice at level one.  We write
$\Delta:=\Delta(R(X,L))$, so that
\begin{equation}
\label{eq:ch33-okounkov-volume}
 \vol(\Delta)=\frac{V}{n!},
 \qquad
 \lim_{m\to\infty}\frac{\dim W_m}{m^n}
 =\vol(\Delta(W_\bullet)),
\end{equation}
the second identity holding for every graded linear series containing an
ample series in the sense of \cite{LM09,KK12}.

\refstepcounter{theorem}\label{thm:ch33-dh-concave-transform}
\begin{quote}
\textbf{Foundational input \thetheorem\ (Boucksom--Chen concave transform;
source boundary).}
Let $\mathcal F$ be a multiplicative, linearly bounded filtration of
$R(X,L)$ with $L$ ample and $X$ integral of dimension $n$.  Define,
for $\tau\in\mathbf R$, the filtered graded linear series
\begin{equation}
\label{eq:ch33-filtered-series}
 W^{(\tau)}_m:=\mathcal F^{m\tau}R_m .
\end{equation}
Then:
\begin{enumerate}
\item each $W^{(\tau)}_\bullet$ is a graded linear subseries of
 $R(X,L)$, and
 $\tau\mapsto\vol\bigl(\Delta(W^{(\tau)}_\bullet)\bigr)$ is decreasing,
 equal to $\vol(\Delta)$ for $\tau\leq C_-$ and to $0$ for
 $\tau>C_+$;
\item there is a bounded, upper semicontinuous, concave function
 $G_{\mathcal F}:\Delta\to\mathbf R$, the \emph{concave transform} of
 $\mathcal F$, such that
\begin{equation}
\label{eq:ch33-concave-transform-levels}
 \vol\bigl(\Delta(W^{(\tau)}_\bullet)\bigr)
 =\vol\{x\in\Delta\mid G_{\mathcal F}(x)\geq\tau\}
\end{equation}
 for all but countably many $\tau$;
\item the empirical measures \eqref{eq:empirical-weight-measure} converge
 weakly, along the sufficiently divisible integers, to
\begin{equation}
\label{eq:ch33-dh-is-pushforward}
 \DHm(\mathcal F)
 =\frac{n!}{V}\,(G_{\mathcal F})_\ast\lambda_\Delta
 =\frac{1}{\vol(\Delta)}\DHm_{\Delta,G_{\mathcal F}} ,
\end{equation}
 in the notation of \Cref{sec:dh-preview}.
\end{enumerate}
Items (1)--(3), including existence and weak convergence, are imported from
\cite[Theorem~A]{BC11}; no proof is claimed here.
\end{quote}

\paragraph{Normalization and mechanism.}
(1)  Multiplicativity \eqref{eq:multiplicative-filtration} gives
$W^{(\tau)}_m\cdot W^{(\tau)}_{m'}\subset W^{(\tau)}_{m+m'}$, since
$m\tau+m'\tau=(m+m')\tau$; so $W^{(\tau)}_\bullet$ is a graded
subalgebra of $R(X,L)$, that is, a graded linear series.  It is decreasing
in $\tau$ because $\mathcal F$ is.  Linear boundedness
\eqref{eq:linear-bounds-filtration} gives
$W^{(\tau)}_m=R_m$ for $\tau\leq C_-$ and $W^{(\tau)}_m=0$ for
$\tau>C_+$.  Existence of the relevant volumes and the treatment of
possibly small filtered subseries are part of the filtered-series theorem
\cite[Theorem~A]{BC11}.  Mere containment in the ample section ring would
not by itself imply the needed volume statement.

(2) and (3)  This is the main theorem of \cite[Theorem~A]{BC11}, whose
statement we now match to ours.  Boucksom and Chen attach to a
multiplicative, linearly bounded filtration of a graded linear series the
concave transform
\begin{equation}
\label{eq:ch33-concave-transform-formula}
 G_{\mathcal F}(x)
 =\sup\Bigl\{\tau\in\mathbf R\ \Bigm|\
    x\in\Delta\bigl(W^{(\tau)}_\bullet\bigr)\Bigr\},
\end{equation}
show that it is concave and bounded on $\Delta$ with the level-set
property \eqref{eq:ch33-concave-transform-levels}, and prove that the
normalized empirical measures of the jumping numbers converge weakly to the
push-forward of normalized Lebesgue measure on $\Delta$ by
$G_{\mathcal F}$.  Concavity of \eqref{eq:ch33-concave-transform-formula}
is the geometric shadow of multiplicativity: if
$x_i\in\Delta(W^{(\tau_i)})$ for $i=0,1$ and $t\in[0,1]$, then
$(1-t)x_0+tx_1\in\Delta(W^{((1-t)\tau_0+t\tau_1)})$, because the product
of sections realizing the two Okounkov data has $\nu$-value the sum and
weight at least the sum.  Upper semicontinuity is part of the construction,
not a later regularization: by \cite[Definition~1.8]{BC11}, the superlevel
set $\{G_{\mathcal F}\geq\tau\}$ is an intersection of Okounkov bodies
$\Delta(W^{(\tau')}_\bullet)$ with $\tau'<\tau$, hence is closed.
Finally, dividing by $\vol(\Delta)=V/n!$ converts the push-forward of
Lebesgue measure into a probability measure, which is
\eqref{eq:ch33-dh-is-pushforward}.
\paragraph{Consequences and normalization checks.}
Existence of the weak limit and its identification as a probability measure
is \Cref{thm:ch33-dh-concave-transform}\,(3).  We give the two elementary
steps separately, since they are used again below.

\emph{Compact support.}  By linear boundedness, every normalized jumping
number $\lambda_{m,j}/m$ lies in the fixed compact interval
$I=[C_-,C_+]$, so every $\nu_m$ is a probability measure supported in
$I$; hence so is any weak limit.  In particular the family
$(\nu_m)$ is tight, and by \Cref{prop:measure-compactness} it has weak
cluster points; \Cref{thm:ch33-dh-concave-transform} identifies the unique
one.

\emph{Tail description.}  For $\tau\in\mathbf R$,
\begin{equation}
\label{eq:ch33-nu-tail}
 \nu_m\bigl([\tau,+\infty)\bigr)
 =\frac{\#\{j\mid\lambda_{m,j}\geq m\tau\}}{N_m}
 =\frac{\dim\mathcal F^{m\tau}R_m}{N_m}
 =\frac{\dim W^{(\tau)}_m}{N_m},
\end{equation}
by \Cref{lem:ch33-filtration-norm-bijection}\,(3).  Dividing numerator and
denominator by $m^n$ and using \eqref{eq:ch33-okounkov-volume} gives
\[
 \lim_{m}\nu_m\bigl([\tau,+\infty)\bigr)
 =\frac{\vol\bigl(\Delta(W^{(\tau)}_\bullet)\bigr)}{\vol(\Delta)} ,
\]
which by \eqref{eq:ch33-concave-transform-levels} and
\Cref{prop:dh-tail} is the tail of the right-hand side of
\eqref{eq:ch33-dh-is-pushforward}, for all but countably many $\tau$.
Convergence of distribution functions at all continuity points, together
with the uniform compact support, upgrades to weak convergence by the
portmanteau theorem \Cref{prop:portmanteau}.

For the second sentence of the statement, let $\mathcal F$ be the
filtration of a relatively ample test configuration with integrally closed
total space.  By \Cref{prop:ch33-central-fibre-weights}\,(3), the weights of
the $\mathbf G_m$-representation $H^0(\cX_0,m\cL_0)$, listed with
multiplicity, are exactly the numbers $\lambda_{m,j}$; hence $\nu_m$ is
the normalized weight distribution of that representation and
$\DHm(\mathcal F)$ is its asymptotic distribution.

For the third sentence, the filtration determines $\lambda_{m,j}$ and
therefore $\nu_m$; conversely, by
\Cref{thm:ch33-tc-filtration-properties}\,(4), equivalent test
configurations, and in particular a test configuration and its integral
closure, give the same filtration.

\begin{corollary}
\label{cor:ch33-dh-structure}
Let $\mathcal F$ be multiplicative and linearly bounded.  Then
$\DHm(\mathcal F)$ is a compactly supported probability measure on
$\mathbf R$ satisfying:
\begin{enumerate}
\item $\DHm(\mathcal F)\bigl((\tau,+\infty)\bigr)
 =\vol(\Delta)^{-1}\vol\{G_{\mathcal F}>\tau\}$, and its atoms are exactly
 the levels $\tau$ with $\vol\{G_{\mathcal F}=\tau\}>0$;
\item $\tau\mapsto\vol\{G_{\mathcal F}\geq\tau\}^{1/n}$ is concave on the
 interval where it is positive; in particular, for $n=1$ the measure
 $\DHm(\mathcal F)$ has a nondecreasing density on the interior of its
 support, away from atoms;
\item $\max\operatorname{supp}\DHm(\mathcal F)=\lambda_{\max}$ of
 \eqref{eq:ch33-fekete-limits}, and
 $\min\operatorname{supp}\DHm(\mathcal F)\geq C_-$ for any admissible
 lower bound $C_-$;
\item for every bounded Borel function $\chi$ on $\mathbf R$,
\begin{equation}
\label{eq:ch33-dh-moment-formula}
 \int_{\mathbf R}\chi\,\rd\DHm(\mathcal F)
 =\frac{1}{\vol(\Delta)}\int_\Delta\chi\bigl(G_{\mathcal F}(x)\bigr)\rd x .
\end{equation}
\end{enumerate}
\end{corollary}

\begin{proof}
All four are the corresponding statements of \Cref{sec:dh-preview} applied
to the pair $(\Delta,G_{\mathcal F})$ and divided by
$\vol(\Delta)$: (1) is \Cref{prop:dh-tail}, (2) is
\Cref{prop:dh-brunn-minkowski}, (4) is \Cref{lem:dh-moments}.  For (3), the
upper end of the support is
$\max_\Delta G_{\mathcal F}$, which equals $\lambda_{\max}$ as follows.  No filtered section has
normalized weight above $\lambda_{\max}$, so the limiting support has
that upper bound.  Conversely, fix $\tau<\lambda_{\max}$ and choose
a section $s\in R_p$ with $\chi_p(s)>p\tau$.  Choose a positive
integer $\ell$ and a positive retained degree $r$ with
\[
 \ell\chi_p(s)+rC_->(\ell p+r)\tau ;
\]
this is possible by making $\ell$ large for a fixed $r$.  Products
$s^{k\ell}R_{kr}$ then lie in $W^{(\tau)}_{k(\ell p+r)}$ for all $k$.
Multiplication is injective, since $X$ is integral, and $rL$ is ample.
This is an ample subseries; its positive volume ensures that the
filtered Okounkov body has positive volume.  Apply the same argument
at a continuity level strictly between $\tau$ and $\lambda_{\max}$
to see that the limiting measure has mass above $\tau$.
Thus its support reaches $\lambda_{\max}$.  This is also
\cite[Lemma~1.6]{BC11}; one isolated finite-level weight would not by
itself prove the assertion.
\end{proof}

\Cref{cor:ch33-dh-structure}\,(2) is the statement announced in
\Cref{rem:dh-log-concavity}: it is the Brunn--Minkowski concavity of the
tail volumes of a concave transform.  It is not a log-concavity assertion
for the density of the measure.

\subsection{Weight counting}
\label{sec:ch33-dh-weights}

\begin{theorem}
\label{thm:ch33-dh-weight-count}
Let $(\cX,\cL)$ be a test configuration for $(X,L)$ with $\cL$
relatively ample and $\cX$ integrally closed, and let
$\mathcal F$ be its filtration.  Write
$N_m=\dim_kH^0(\cX_0,m\cL_0)=h^0(X,mL)$ and let $w_m$ be the total
weight \eqref{eq:ch33-total-weight} of the $\mathbf G_m$-representation
$H^0(\cX_0,m\cL_0)$.  Then, for $m$ sufficiently divisible and large:
\begin{enumerate}
\item $\displaystyle
 \int_{\mathbf R}\lambda\,\rd\nu_m(\lambda)=\frac{w_m}{mN_m}$;
\item $N_m=\dfrac{V}{n!}m^n+O(m^{n-1})$ and
 $w_m=a_0m^{n+1}+O(m^n)$ for a rational number $a_0$, and
\begin{equation}
\label{eq:ch33-first-moment-limit}
 \int_{\mathbf R}\lambda\,\rd\DHm(\mathcal F)(\lambda)
 =\lim_{m\to\infty}\frac{w_m}{mN_m}
 =\frac{n!\,a_0}{V};
\end{equation}
\item with the unnormalized convention
 \eqref{eq:ch33-dh-unnormalized},
\begin{equation}
\label{eq:ch33-first-moment-intersection}
 \int_{\mathbf R}\lambda\,\rd\DHm^{\mathrm u}(\mathcal F)(\lambda)
 =\frac{(\overline\cL^{\,n+1})}{n+1},
\end{equation}
 and more generally the $q$-th moment of $\DHm(\mathcal F)$ is the
 normalized intersection number of
 \cite[Corollary~3.4]{BHJ17}, as recorded in
 \eqref{eq:cma-bhj-moments}.
\end{enumerate}
\end{theorem}

\begin{proof}
(1)  By \Cref{prop:ch33-central-fibre-weights}\,(3) the weights of
$H^0(\cX_0,m\cL_0)$ with multiplicity are $\lambda_{m,1},\dots,
\lambda_{m,N_m}$; hence
$\int\lambda\,\rd\nu_m=N_m^{-1}\sum_j\lambda_{m,j}/m=w_m/(mN_m)$.

(2)  The asymptotics of $N_m$ are Hilbert--Samuel asymptotics for the
ample $L$, giving \eqref{eq:ch33-volume}.  For $w_m$, the total weight
of the representation on $H^0(\cX_0,m\cL_0)$ is a polynomial in $m$ of
degree at most $n+1$ for $m\gg0$; this is the equivariant
Riemann--Roch calculation of
\cite[Proposition~3.12(iii), proof]{BHJ17}, recalled below.  Since all
normalized weights lie in the fixed compact interval $I$, the function
$\lambda\mapsto\lambda$ is bounded on the supports of all $\nu_m$ and
of the limit, so weak convergence gives convergence of first moments; this
is one use of uniform compact support, and it is exactly
what \Cref{rem:dh-moments-not-continuous} shows cannot be dispensed with in
general.

(3)  The direct source is
\cite[Proposition~3.12(iii), equation~(3.3)]{BHJ17}.  To recall its
algebraic argument, set
$\mathcal E_m=\overline\pi_\ast(m\overline\cL)$.
For large divisible $m$, relative Serre vanishing makes this a vector
bundle of rank $N_m$ with central fibre $H^0(\cX_0,m\cL_0)$.
The fixed product trivialization near infinity gives
\[
 w_m=\deg\det\mathcal E_m
     =\chi(\mathbf P^1,\mathcal E_m)-N_m
     =\chi(\overline\cX,m\overline\cL)-N_m .
\]
The first equality is the weight-degree identity for the equivariant
bundle, with the convention of \eqref{eq:ch33-h0-weight-decomposition};
the second is Riemann--Roch on $\mathbf P^1$, and the third is Leray
with higher direct images zero.  Asymptotic Riemann--Roch on the
projective $(n+1)$-dimensional scheme then gives
\[
 w_m=\frac{(\overline\cL^{\,n+1})}{(n+1)!}\,m^{n+1}+O(m^n).
\]
Here $n\geq1$, so subtracting $N_m=O(m^n)$ does not change the
leading coefficient.  Combine with (2) and multiply by $V$ to obtain
\eqref{eq:ch33-first-moment-intersection}.

Higher moments use instead the auxiliary BHJ fibrations over
$\mathbf P^q$ associated to the polarized central fibre, as in
\cite[Corollary~3.4]{BHJ17} and \eqref{eq:cma-bhj-moments}.
They are not computed by replacing $n+1$ with $n+q$ on the original
compactification.  Both algebraic results are independent of the later
energy-asymptotics erratum discussed in \Cref{thm:rf-erratum-audit}.

\end{proof}

\begin{remark}
\label{rem:ch33-dh-unit-tests}
Formula \eqref{eq:ch33-first-moment-intersection} is compatible with the
twist and base-change rules of
\Cref{thm:ch33-intersection-invariance}\,(3),(4) and
\Cref{prop:ch33-dh-transformations}: twisting by $c$ adds $c$ to the
first moment of $\DHm$ and $(n+1)cV$ to $(\overline\cL^{\,n+1})$,
hence $cV$ to the left-hand side of
\eqref{eq:ch33-first-moment-intersection} after multiplication by $V$;
degree-$d$ base change multiplies both sides by $d$.  Both tests are
carried out numerically in \Cref{ex:ch33-dnc-p1}.
\end{remark}

\subsection{Behaviour under the basic operations}
\label{sec:ch33-dh-operations}

Adding $cm$ to all weights translates $\DHm(\mathcal F)$ by $c$;
base change of degree $d$ dilates it by $d$.  Thus one must fix both the
linearization and the normalization before comparing moments.

\begin{proposition}
\label{prop:ch33-dh-transformations}
Let $\mathcal F$ be multiplicative and linearly bounded, and write
$T_c(\lambda)=\lambda+c$, $m_d(\lambda)=d\lambda$.  Then
\begin{equation}
\label{eq:ch33-dh-transform}
\begin{gathered}
 \DHm(\mathcal F^{(c)})=(T_c)_\ast\DHm(\mathcal F),
 \qquad
 \DHm(\mathcal F^{[d]})=(m_d)_\ast\DHm(\mathcal F),\\
 \DHm(\mathcal F^{(r)}_{\mathrm{Ver}})
   =(m_r)_\ast\DHm(\mathcal F).
\end{gathered}
\end{equation}
where $\mathcal F^{(c)}$, $\mathcal F^{[d]}$ and
$\mathcal F^{(r)}_{\mathrm{Ver}}$ are the twist, base change, and
Veronese restriction of \Cref{prop:ch33-transformations}.  All three
operations preserve total mass, and the first moment transforms by
$\mu\mapsto\mu+c$, $\mu\mapsto d\mu$, $\mu\mapsto r\mu$,
respectively.  If the retained Veronese weights are divided by $r$, or
if they are normalized by the actual degree $rm$ rather than the new
degree $m$, the measure and its first moment are unchanged.
\end{proposition}

\begin{proof}
The jumping numbers of $\mathcal F^{(c)}$ in degree $m$ are $\lambda_{m,j}+mc$, so the normalized ones are $\lambda_{m,j}/m+c$ and $\nu_m$ is replaced by $(T_c)_*\nu_m$. Pass to the limit. The algebraic operation $\chi^{[d]}=d\chi$ likewise dilates every finite-level measure by $d$.  For the reindexed Veronese filtration on $R(X,rL)$, the degree-$m$ empirical measure is
\[
 \frac1{N_{rm}}\sum_j\delta_{\lambda_{rm,j}/m}
 =(m_r)_*\nu_{rm}.
\]
The subsequence $\nu_{rm}$ has the same limit as $\nu_m$, but reindexing
changes the denominator from $rm$ to $m$; this produces the factor $r$.
Dividing the retained weights by $r$ cancels it.  Uniformly compact
supports permit passage to first moments, giving the stated transformations
and preservation of mass.

For a test configuration, base change followed by integral closure in the
generic fibre has the same limiting DH measure as the raw base change by
\cite[Corollary~5.8 and Proposition~5.9]{BHJ17}.  Indeed, compare both
families through their common full normalization and use normalization
invariance.  Hence the DH measure also dilates by $d$.  This assertion
concerns the limiting measure; the finite-level weights themselves may
change, as in \Cref{ex:rf-normalization-changes-weights}.
\end{proof}

\begin{proposition}
\label{prop:ch33-dh-continuity}
For multiplicative linearly bounded graded norms $\chi,\chi'$ on the
ample section ring $R(X,L)$, set
\begin{equation}
\label{eq:ch33-dinfty}
 d_\infty(\chi,\chi')
 :=\sup_{m\geq1}\ \sup_{0\ne s\in R_m}
        \frac{|\chi_m(s)-\chi'_m(s)|}{m}.
\end{equation}
If $d_\infty(\chi,\chi')\leq\varepsilon$, then for every $m$ and every
$j$,
\begin{equation}
\label{eq:ch33-jumping-comparison}
 |\lambda_{m,j}-\lambda'_{m,j}|\leq\varepsilon m ,
\end{equation}
and consequently, for every $1$-Lipschitz function $f$ on
$\mathbf R$,
\begin{equation}
\label{eq:ch33-dh-lipschitz}
 \Bigl|\int f\,\rd\DHm(\chi)-\int f\,\rd\DHm(\chi')\Bigr|
 \leq\varepsilon .
\end{equation}
In particular $\chi\mapsto\DHm(\chi)$ is $1$-Lipschitz for
$d_\infty$ and the Kantorovich--Rubinstein distance, hence continuous for
the weak topology; and if $d_\infty(\chi_j,\chi)\to0$ with all
$\chi_j$ linearly bounded by the same constants, then all moments of
$\DHm(\chi_j)$ converge to those of $\DHm(\chi)$.
\end{proposition}

\begin{proof}
Write $\delta=\varepsilon m$.  The hypothesis gives, for every
$\lambda$,
\[
 \{s\ne0\mid\chi'_m(s)\geq\lambda\}
 \subset\{s\ne0\mid\chi_m(s)\geq\lambda-\delta\},
\]
whence $\dim\mathcal F'^\lambda R_m\leq\dim\mathcal F^{\lambda-\delta}R_m$,
and symmetrically.  Since
$\lambda_{m,j}=\sup\{\lambda\mid\dim\mathcal F^\lambda R_m\geq j\}$, these
two inequalities give
$\lambda'_{m,j}\leq\lambda_{m,j}+\delta$ and
$\lambda_{m,j}\leq\lambda'_{m,j}+\delta$, which is
\eqref{eq:ch33-jumping-comparison}.  For a $1$-Lipschitz $f$,
\[
 \Bigl|\int f\,\rd\nu_m-\int f\,\rd\nu'_m\Bigr|
 \leq\frac1{N_m}\sum_{j=1}^{N_m}
      \Bigl|f\Bigl(\frac{\lambda_{m,j}}m\Bigr)
            -f\Bigl(\frac{\lambda'_{m,j}}m\Bigr)\Bigr|
 \leq\frac{\delta}{m}=\varepsilon,
\]
using the matching of the two decreasing lists.  Letting $m\to\infty$
gives \eqref{eq:ch33-dh-lipschitz}: each pair of measures has uniformly
compact support, on which the possibly unbounded Lipschitz function
$f$ agrees with a bounded continuous extension.  That extension is a
legitimate test function for weak convergence.  The final assertion follows
because all the measures are then supported in one compact interval, on
which every polynomial is Lipschitz.
\end{proof}

\begin{warning}{No continuity for pointwise convergence}
\label{warn:ch33-dh-no-pointwise-continuity}
The uniform metric in \eqref{eq:ch33-dinfty} cannot be replaced by
pointwise convergence.  Here are two different failures on
$(X,L)=(\mathbf P^1,\cO(1))$.

First, give every nonzero section in degree $m$ the same weight
$\chi_m^{(j)}=-\min\{m,j\}$, with degree-zero weight zero.
These are multiplicative norms, since
$\min\{m+m',j\}\leq\min\{m,j\}+\min\{m',j\}$, and share linear
bounds.  For each fixed section they converge to $\chi_m=-m$.
For fixed $j$, however, $\chi_m^{(j)}/m\to0$ as $m\to\infty$,
so their DH measures are $\delta_0$, while the limit norm has DH
measure $\delta_{-1}$.

Second, choose pairwise distinct closed points $p_j$.  The semiample
boundary case of \Cref{ex:ch33-dnc-p1} gives functions
$\phi_j(v)=-\min\{v(\mathfrak m_{p_j}),1\}$.
Each is nonconstant only on the ray through $p_j$, so
$\phi_j(v)\to0$ for every fixed point of $(\mathbf P^1)^{\An}$.
Every one of their DH measures is uniform on $[-1,0]$, whereas
the zero function has measure $\delta_0$.  Thus pointwise convergence
of the analytic functions also fails to control DH.

Weak convergence alone does not control moments without uniform
linear bounds; \Cref{rem:dh-moments-not-continuous} gives a separate
example.  The finer distances $d_p$ belong to
\Cref{chap:bj-stability}.

\end{warning}

\section{Worked examples}
\label{sec:ch33-examples}

Every construction of this chapter is carried out to explicit numbers in the
four examples below.  The first two are test configurations on
$\mathbf P^1$, presented simultaneously as a total space, a flag ideal, a
Rees algebra, a filtration, a Fubini--Study function, and a
Duistermaat--Heckman measure; each of them is checked against the
intersection formula \eqref{eq:ch33-first-moment-intersection}.  The third is
a filtration that is not of finite type, with a complete proof of that
failure.  The fourth computes the whole dictionary on a toric variety
directly from the polytope.

\subsection{A product configuration}
\label{sec:ch33-example-product}

\begin{example}
\label{ex:p1-product-dh}
Let $X=\mathbf P^1$, $L=\cO(1)$, and let
$\mathbf G_m$ act by $a\cdot[x:y]=[ax:y]$.  Choose the linearization for
which $x^jy^{m-j}$ has weight $j$.  Then
\[
 \nu_m=\frac1{m+1}\sum_{j=0}^m\delta_{j/m}
       \longrightarrow \mathbf 1_{[0,1]}\,\rd \lambda .
\]
Although the family is a product as a scheme, its nontrivial linearization
produces a nonconstant norm and a nontrivial Duistermaat--Heckman measure.

We now record all the associated data, in order to have one completely
computed model of the dictionary.  Here $R_m=H^0(\mathbf P^1,\cO(m))$ has
the monomial basis $x^jy^{m-j}$, $0\leq j\leq m$, and $N_m=m+1$,
$V=(L)=1$, $n=1$.

\emph{The filtration.}  By hypothesis $\chi_m(x^jy^{m-j})=j$, and since
the configuration is equivariant for the torus fixing $[0:1]$ and
$[1:0]$, the filtration is monomial:
\begin{equation}
\label{eq:ch33-p1-product-filtration}
 \mathcal F^\lambda R_m
 =\operatorname{span}\{x^jy^{m-j}\mid j\geq\lambda\},
 \qquad
 \chi_m\Bigl(\sum_j c_jx^jy^{m-j}\Bigr)=\min\{j\mid c_j\ne0\}.
\end{equation}
Multiplicativity is the identity $j+j'=j+j'$ together with the fact that
the support of a product of two polynomials in one variable is contained in
the sum of the supports; linear boundedness holds with $C_-=0$,
and any $C_+>1$, for instance $C_+=2$.  The endpoint $C_+=1$ is not
admissible in \eqref{eq:linear-bounds-filtration}, since $x^m$ has weight
$m$ and remains in $\mathcal F^mR_m$.

\emph{The Rees algebra.}  By \eqref{eq:filtration-rees}, $\mathcal
R(\mathcal F)$ is spanned by the monomials
$x^jy^{m-j}\varpi^{-\lambda}$ with $\lambda\leq j$, so
\begin{equation}
\label{eq:ch33-p1-product-rees}
 \mathcal R(\mathcal F)=k[\,x\varpi^{-1},\ y,\ \varpi\,],
\end{equation}
a polynomial ring in three variables:
$x^jy^{m-j}\varpi^{-\lambda}
=(x\varpi^{-1})^{j}y^{m-j}\varpi^{\,j-\lambda}$ with $j-\lambda\geq0$.
It is generated in degree one over $k[\varpi]$, so no Veronese change is
needed, and
$\Proj_{\mathbf A^1}\mathcal R(\mathcal F)=\mathbf P^1\times\mathbf A^1$.
As a scheme this is the product; as a test configuration it is not the
trivial one, because the trivialization
\eqref{eq:tc-punctured-trivialization} is
$[x':y]=[\varpi^{-1}x:y]$, not the identity.  This is
\Cref{ex:ch33-product-nontrivial-action} in coordinates.

\emph{The Fubini--Study function.}  By \Cref{cor:ch33-fg-finite-max} with
$r=1$ and the adapted basis $x,y$,
\begin{equation}
\label{eq:ch33-p1-product-fs}
 \phi_\chi(v)=\max\{\log|x|(v)+1,\ \log|y|(v)\}.
\end{equation}
On the ray $t\ord_{[0:1]}$, where $v(x)=t$ and $v(y)=0$, this is
$\max\{1-t,0\}$; on the ray $t\ord_{[1:0]}$ it is
$\max\{1,-t\}=1$; at every other point of $(\mathbf P^1)^{\An}$ it is
$1$.  Thus $\phi_\chi$ is continuous, equal to $1$ away from one ray,
and equal to $1+\phi_{\mathfrak a}$ for the flag ideal
$\mathfrak a=\mathfrak m_{[0:1]}+(\varpi)$, whose function is
$\max\{\log|\mathfrak m_{[0:1]}|,-1\}$.

\emph{The Duistermaat--Heckman measure and its checks.}  The jumping numbers
in degree $m$ are $0,1,\dots,m$, each once, so the displayed limit is
correct and
\begin{equation}
\label{eq:ch33-p1-product-dh}
 \DHm(\mathcal F)=\mathbf 1_{[0,1]}\,\rd\lambda,
 \qquad
 \int\lambda\,\rd\DHm(\mathcal F)=\tfrac12 .
\end{equation}
Since $V=1$, the probability and unnormalized measures coincide, and
\eqref{eq:ch33-first-moment-intersection} predicts
$(\overline\cL^{2})=2\cdot\tfrac12=1$.  This is confirmed in two
independent ways.  First, by the semiample $c=1$ boundary case of
\Cref{ex:ch33-dnc-p1} below, the configuration
with Fubini--Study function $\phi_{\mathfrak a}$ has
$(\overline\cL^2)=-1$, and twisting by the character of weight $1$ adds
$(n+1)cV=2$ by \eqref{eq:ch33-intersection-twist}, giving $-1+2=1$.
Second, the compactification $\overline\cX$ is the Hirzebruch surface
$\mathbf P(\cO\oplus\cO(-1))$ over $\mathbf P^1$: the two charts are
glued by $[x':y]=[\varpi^{-1}x:y]$, a transition cocycle of degree
$-1$.  Writing $\xi$ for the tautological class and $f$ for the fibre
class, one has $\xi^2=-1$, $\xi\cdot f=1$, $f^2=0$, and the
normalization above gives $\overline\cL=\xi+f$, so
$(\overline\cL^2)=-1+2=1$.

\emph{Concave transform.}  Taking $\nu=\ord_{[0:1]}$ as the flag
valuation, the Okounkov body of $(\mathbf P^1,\cO(1))$ is
$\Delta=[0,1]$ with $\vol(\Delta)=1=V/1!$, and the concave transform of
\eqref{eq:ch33-p1-product-filtration} is
$G_{\mathcal F}(u)=u$.  Then
\eqref{eq:ch33-dh-is-pushforward} reads
$\DHm(\mathcal F)=(\operatorname{id})_\ast\lambda_{[0,1]}$, which is
\eqref{eq:ch33-p1-product-dh}.

\emph{The reflection test.}  Replacing the one-parameter subgroup by its
inverse reverses the weights of each monomial in the adapted basis:
the new norm is $\min\{-j:c_j\ne0\}$, rather than the negative of
the old minimum.  The reflected weight list reflects
$\DHm$, and changes the first moment from $1/2$ to $-1/2$; by
\eqref{eq:ch33-first-moment-intersection} it changes
$(\overline\cL^2)$ from $1$ to $-1$.  This is the unit test announced
in \Cref{rem:ch33-weight-sign}.
\end{example}

\subsection{A nontrivial configuration on the projective line}
\label{sec:ch33-example-dnc}

\begin{example}
\label{ex:ch33-dnc-p1}
Let $X=\mathbf P^1$, $L=\cO(1)$, $p=[0:1]$, and let $z=x/y$ be the
affine coordinate near $p$.  Fix $c\in\mathbf Q$ with $0<c<1$.

\emph{The flag ideal and its Rees algebra.}  Let
\begin{equation}
\label{eq:ch33-dnc-flag}
 \mathfrak a:=\mathfrak m_p+(\varpi)\subset\cO_{X_{\triv}} ,
\end{equation}
which in the local coordinates at $(p,0)$ is the maximal ideal
$(z,\varpi)$ and is the unit ideal elsewhere.  In the indexing
\eqref{eq:ch33-flag-weight} its weight decomposition is
$\mathfrak a_0=\mathfrak m_p$ and $\mathfrak a_{-1}=\cO_X$, so
\begin{equation}
\label{eq:ch33-dnc-phi-a}
 \phi_{\mathfrak a}=\max\{\log|\mathfrak m_p|,\,-1\},
 \qquad
 \phi_{\mathfrak a}(t\ord_p)=\max\{-t,-1\}=-\min\{t,1\},
\end{equation}
and $\phi_{\mathfrak a}=0$ at every point of
$(\mathbf P^1)^{\An}$ not on the ray through $p$.  Every power
$\mathfrak a^q$ is integrally closed, being locally the $q$-th power of
the maximal ideal of a two-dimensional regular local ring, so the normalized
Rees algebra \eqref{eq:flag-rees-algebra} is the ordinary one:
\begin{equation}
\label{eq:ch33-dnc-rees}
 \mathcal R(\mathfrak a)=\bigoplus_{q\geq0}(\mathfrak m_p+(\varpi))^q u^q .
\end{equation}

\emph{The total space.}  By
\Cref{thm:ch33-flag-tc-equivalence}\,(1), $\cX:=\Proj_{X_{\triv}}\mathcal
R(\mathfrak a)$ is the blowup of $X_{\triv}=\mathbf P^1\times\mathbf A^1$
at the point $(p,0)$, an integrally closed---indeed smooth---test
configuration dominating the trivial one.  This is the deformation to the
normal cone of \Cref{ex:product-dnc-test-configurations}.  Its central
fibre is $\cX_0=\widetilde E+E$, where $E$ is the exceptional
$(-1)$-curve and $\widetilde E$ is the strict transform of
$\mathbf P^1\times\{0\}$, which is also a $(-1)$-curve since
$\cX_0^2=0$, $\widetilde E\cdot E=1$.  Both $b_{E}$ and
$b_{\widetilde E}$ of \eqref{eq:ch33-bE} equal $1$.

\emph{The polarization.}  Set
$\cL_c:=\mu^\ast p_1^\ast\cO(1)-cE$.  Then
\[
 \cL_c\cdot E=c>0,
 \qquad
 \cL_c\cdot\widetilde E=1-c>0 ,
\]
so $\cL_c$ is relatively ample for $0<c<1$, semiample for $c=1$.  Its
Fubini--Study function is $\phi_{\cL_c}=c\,\phi_{\mathfrak a}$, by
\Cref{thm:ch33-flag-tc-equivalence}\,(1) and linearity, so
\begin{equation}
\label{eq:ch33-dnc-phi}
 \phi_{\cL_c}(t\ord_p)=-c\min\{t,1\},
 \qquad
 \phi_{\cL_c}=0\ \text{elsewhere}.
\end{equation}

\emph{The filtration.}  A section $x^ay^{m-a}$ of $mL$ has
$\operatorname{div}=a\,[p]+(m-a)\,[\infty]$, so on $\cX$
\[
 \operatorname{div}_{\mu^\ast(mL_{\triv})}(\bar s)
 =aH_0+(m-a)H_\infty+aE ,
\]
where $H_0,H_\infty$ are the strict transforms of the two horizontal
sections; the coefficient $a$ on $E$ appears because
$\ord_E(z)=1$.  Viewing $\bar s$ as a section of
$m\cL_c=\mu^\ast(mL_{\triv})-mcE$ subtracts $mcE$, and
$\operatorname{div}(\varpi)=\widetilde E+E$ subtracts $\lambda$ from
both central components.  Hence
$\varpi^{-\lambda}\bar s$ is regular if and only if
$\lambda\leq0$ and $\lambda\leq a-mc$, that is,
\begin{equation}
\label{eq:ch33-dnc-weights}
 \chi_m(x^ay^{m-a})=\min\{0,\ a-mc\} .
\end{equation}
Again the filtration is monomial, by torus equivariance of the whole
construction.  For $c=1/2$ and $m=2$ the weights are
$\chi_2(x^2)=\chi_2(xy)=0$ and $\chi_2(y^2)=-1$.

\emph{Base ideals recover the flag ideal.}  For $c=1/2$ and $m=2$:
$\mathcal F^{-1}R_2=R_2$, so $\mathfrak b_{2,-1}=\cO_X$;
$\mathcal F^0R_2=\langle x^2,xy\rangle$, whose base ideal is
$(z^2,z)=\mathfrak m_p$; and $\mathcal F^\lambda R_2=0$ for
$\lambda\geq1$.  Therefore
$\mathfrak a_{(2)}=\mathfrak m_p+\cO_X\varpi=\mathfrak a$, confirming
\Cref{prop:ch33-base-ideal-flag}\,(2) and
\eqref{eq:ch33-flag-vs-norm}:
$\tfrac12\phi_{\mathfrak a}=\phi_{\cL_{1/2}}$.  Here the generating
Veronese index is $r=2$; for every $\ell\geq1$, the same monomial
calculation gives
$\overline{\mathfrak a_{(2\ell)}}=\mathfrak a^\ell$, illustrating
\eqref{eq:ch33-flag-power-propagation}.  In particular
$\mathfrak a_{(2)}=\mathfrak a$, not $\mathfrak a^2$.

\emph{The Rees algebra of the filtration.}  Still for $c=1/2$, the
Veronese index $r=2$ works and
\begin{equation}
\label{eq:ch33-dnc-filtration-rees}
 \mathcal R(\mathcal F)^{(2)}
 =k[\,x^2,\ xy,\ y^2\varpi,\ \varpi\,]
 \;\simeq\;
 k[U,W,T,\varpi]/(UT-\varpi W^2).
\end{equation}
Indeed a monomial $x^ay^{2m-a}\varpi^{-\lambda}$ lies in
$\mathcal R(\mathcal F)^{(2)}$ exactly when
$\lambda\leq\min\{0,a-m\}$; putting $k'=\max\{0,m-a\}$,
$i=\max\{a-m,0\}$, $j=2m-a-2k'$, one has $i+j+k'=m$, $2i+j=a$ and
\[
 x^ay^{2m-a}\varpi^{-\lambda}
 =(x^2)^{i}(xy)^{j}(y^2\varpi)^{k'}\varpi^{-\lambda-k'},
 \qquad-\lambda-k'\geq0 .
\]
The relation $UT=\varpi W^2$ is $x^2\cdot y^2\varpi=\varpi(xy)^2$.  The
relative $\Proj$ is the family of conics
$\{UT=\varpi W^2\}\subset\mathbf P^2\times\mathbf A^1$, smooth as a
surface, with smooth conics over $\varpi\ne0$ and the union of the two
lines $\{U=0\}$, $\{T=0\}$ over $\varpi=0$: exactly
$\widetilde E\cup E$.

\emph{The Duistermaat--Heckman measure.}  By
\eqref{eq:ch33-dnc-weights} the normalized jumping numbers in degree $m$
are $\min\{0,a/m-c\}$ for $0\leq a\leq m$, each occurring once, so
\begin{equation}
\label{eq:ch33-dnc-dh}
 \DHm(\mathcal F)
 =(1-c)\,\delta_0+\mathbf 1_{[-c,0]}\,\rd\lambda ,
\end{equation}
a probability measure, in agreement with
\Cref{cor:ch33-dh-structure}: the concave transform on
$\Delta=[0,1]$ is $G_{\mathcal F}(u)=\min\{0,u-c\}$, whose level set
$\{G_{\mathcal F}=0\}$ has length $1-c$, which is the atom.

\emph{The two numerical checks.}  The first moment is
\[
 \int\lambda\,\rd\DHm(\mathcal F)
 =\int_0^1\min\{0,u-c\}\,\rd u
 =\int_0^c(u-c)\,\rd u=-\frac{c^2}{2}.
\]
On the other side, $\overline\cX$ is the blowup of
$\mathbf P^1\times\mathbf P^1$ at $(p,0)$ and
$\overline\cL_c=\mu^\ast p_1^\ast\cO(1)-cE$, so
\[
 (\overline\cL_c^2)
 =(p_1^\ast\cO(1))^2-2c\,(\mu^\ast p_1^\ast\cO(1)\cdot E)+c^2(E^2)
 =0-0-c^2=-c^2 ,
\]
and $(\overline\cL_c^2)/((n+1)V)=-c^2/2$, confirming
\eqref{eq:ch33-first-moment-intersection} with $V=1$.  Twisting the
linearization by $c'$ adds $c'$ to the first moment and
$2c'$ to $(\overline\cL_c^2)$, by
\eqref{eq:ch33-intersection-twist}; base change of degree $d$ multiplies
both by $d$, by \eqref{eq:ch33-intersection-base-change}.  This is the
unit test of \Cref{rem:ch33-dh-unit-tests}.

\emph{The boundary case.}  For $c=1$ the polarization is only semiample:
it contracts $\widetilde E$, and the resulting configuration is the
product configuration of \Cref{ex:p1-product-dh} twisted by the character of
weight $-1$.  Indeed \eqref{eq:ch33-dnc-weights} with $c=1$ gives
$\chi_m(x^ay^{m-a})=a-m$, which is
\eqref{eq:ch33-p1-product-filtration} shifted by $-m$, and
\eqref{eq:ch33-dnc-dh} becomes $\mathbf 1_{[-1,0]}\rd\lambda$, which is
\eqref{eq:ch33-p1-product-dh} translated by $-1$.
\end{example}

\subsection{A filtration that is not of finite type}
\label{sec:ch33-example-not-fg}

Before the example we record the notion that fails, and the criterion used
to detect the failure.

\begin{definition}
\label{def:ch33-finite-type}
A multiplicative filtration $\mathcal F$ of $R(X,L)$ is of
\emph{finite type} if there are $r\geq1$ and finitely many nonzero
sections $s_1,\dots,s_N\in R_r$ such that, for every $m\geq1$ and every
$\lambda$, the space $\mathcal F^\lambda R_{rm}$ is spanned by the
products $s_{i_1}\cdots s_{i_m}$ with
$\sum_{j}\chi_r(s_{i_j})\geq\lambda$.
\end{definition}

By \Cref{thm:ch33-fg-filtration-tc}, a rational finitely generated
filtration is of finite type, and a filtration of finite type with integer
weights is rational and finitely generated.  The point of
\Cref{def:ch33-finite-type} is that it also makes sense for real weights,
where it is the notion of an $\mathbf R$-test configuration.

\begin{lemma}\label{lem:ch33-finite-type-pl}
If $\mathcal F$ is of finite type, then
\begin{equation}
\label{eq:ch33-finite-type-max}
 \phi_\chi(v)=\frac1r\max_{1\leq i\leq N}\{\chi_r(s_i)-v(s_i)\}
\end{equation}
for all $v\in X^{\An}$.  In particular, along any divisorial ray
$t\mapsto\phi_\chi(t\,\ord_F)$ the function $\phi_\chi$ is a finite
maximum of affine functions of $t$, hence convex and piecewise affine with
at most $N$ pieces.
\end{lemma}

\begin{proof}
The proof of \Cref{cor:ch33-fg-finite-max} used only the spanning property
of \Cref{def:ch33-finite-type}, and it applies verbatim.  For the last
assertion, $\ord_F(s_i)$ is a fixed nonnegative number for each $i$, so
$t\mapsto\chi_r(s_i)-t\ord_F(s_i)$ is affine.
\end{proof}

\begin{example}
\label{ex:ch33-harmonic-not-fg}
Let $X=\mathbf P^1$, $L=\cO(1)$, and $R_m=H^0(X,\cO(m))$ with basis
$x^ay^b$, $a+b=m$.  Define
\begin{equation}
\label{eq:ch33-harmonic-weights}
 \Psi(a,b):=\frac{ab}{a+b}\quad\bigl((a,b)\ne(0,0)\bigr),
 \qquad \Psi(0,0):=0 ,
\end{equation}
and set
\begin{equation}
\label{eq:ch33-harmonic-filtration}
\begin{gathered}
 \chi_m\Bigl(\sum_{a+b=m}c_{a}x^ay^b\Bigr)
 :=\min\{\Psi(a,b)\mid c_a\ne0\},\\
 \mathcal F^\lambda R_m
 =\operatorname{span}\{x^ay^b\mid\Psi(a,b)\geq\lambda\}.
\end{gathered}
\end{equation}

\emph{This is a legitimate multiplicative, linearly bounded filtration.}
The function $\Psi$ is, up to the factor $2$, the harmonic mean, hence
concave on $\mathbf R^2_{\geq0}$, and it is positively homogeneous of
degree one.  A concave, positively homogeneous function on a convex cone is
superadditive, so
\[
 \Psi(a+a',b+b')\geq\Psi(a,b)+\Psi(a',b') .
\]
Given $s\in R_m$ and $s'\in R_{m'}$, every monomial occurring in
$ss'$ is a product of a monomial of $s$ and a monomial of $s'$---the
coefficient of $x^{\ell}y^{m+m'-\ell}$ in $ss'$ is a sum of products
$c_ac'_{a'}$ with $a+a'=\ell$, and it is nonzero only if some summand
is---so
$\chi_{m+m'}(ss')\geq\chi_m(s)+\chi_{m'}(s')$.  Since
$0\leq\Psi(a,b)\leq(a+b)/4$, the filtration is linearly bounded with
$C_-=0$ and any $C_+>1/4$, say $C_+=1/2$.  The value $C_+=1/4$ is
\emph{not} admissible in \eqref{eq:linear-bounds-filtration}, the bound
$\Psi\leq(a+b)/4$ being attained: for even $m$ the monomial
$x^{m/2}y^{m/2}$ has $\Psi=m/4$ and so lies in
$\mathcal F^{m/4}R_m$.  The sharp spectral endpoint remains
$\lambda_{\max}=1/4$, and $\operatorname{supp}\DHm=[0,1/4]$.
Note that every weight
$\Psi(a,m-a)=a(m-a)/m$ is a rational number.

\emph{Its Fubini--Study function.}  As in the previous examples the
filtration is monomial, so
\[
 \phi_\chi(t\ord_p)
 =\sup_{m\geq1}\max_{0\leq a\leq m}
    \frac{a(m-a)/m-ta}{m}
 =\sup_{u\in[0,1]}\bigl\{u(1-u)-tu\bigr\},
\]
where $p=[0:1]$ and we used $v(x^ay^{m-a})=ta$ in the frame $y^m$; the
supremum over all sections reduces to the supremum over monomials because
$\chi_m$ is a minimum over the monomial support and
$v(s)\geq\min v(x^ay^b)$.  This inequality bounds every section's
contribution by the largest monomial contribution; monomials themselves
attain the reverse bound.  On this ray the valuation minimum is also exact.
The quadratic $u\mapsto u(1-t-u)$ has its maximum on $[0,1]$ at
$u=(1-t)/2$ when $0\leq t\leq1$, and at $u=0$ when $t\geq1$.  Hence
\begin{equation}
\label{eq:ch33-harmonic-fs}
 \phi_\chi(t\ord_p)=
 \begin{cases}
  \dfrac{(1-t)^2}{4}, & 0\leq t\leq1,\\[2mm]
  0,& t\geq1,
 \end{cases}
\end{equation}
and by symmetry the same formula holds on the ray through $[1:0]$; on
every other ray, and at the trivial valuation, $\phi_\chi=1/4$.  The
function \eqref{eq:ch33-harmonic-fs} is continuously differentiable,
convex, and decreasing, as it must be for a psh function on a curve, and it
takes the value $1/4$ at $t=0$.

\emph{It is not of finite type.}  Suppose it were.  By
\Cref{lem:ch33-finite-type-pl}, $t\mapsto\phi_\chi(t\ord_p)$ would be
piecewise affine on $[0,\infty)$.  But \eqref{eq:ch33-harmonic-fs} is
strictly convex on $[0,1]$, hence affine on no subinterval.  This
contradiction proves that $\mathcal F$ is not of finite type; in particular,
it cannot be the filtration of a semiample test configuration.

\emph{What exactly fails.}  Three separate things, and it is worth keeping
them apart.
\begin{enumerate}
\item \emph{No finite generating set.}  In the semigroup language of
 \Cref{prop:ch33-monomial-fg-criterion}, the relevant object is the
 hypograph of $\Psi$ inside $\mathbf R^2_{\geq0}\times\mathbf R$.  It is
 a closed convex cone, but it is not polyhedral: its boundary contains the
 nonaffine, strictly concave profile
 $\Psi(u,1-u)=u(1-u)$ on the degree-one slice.
 The integer-level hypograph semigroup of
 \Cref{prop:ch33-monomial-fg-criterion} is therefore not finitely
 generated either.  This is an additional convex-geometric obstruction;
 it does not identify integer-rounded weights with the original
 rational weights.
\item \emph{Unbounded denominators, not irrational weights.}  Every jumping
 number $a(m-a)/m$ in degree $m$ is rational, and the normalized ones
 $a(m-a)/m^2$ are rational as well.  What fails is that no single $r$
 makes all weights lie in $r^{-1}\mathbf Z$: in degree $m$ the
 denominators are of size $m$.  This is a strictly weaker failure than
 irrationality, and it already destroys finite generation.
 \item \emph{The remaining properties persist.}  The filtration is multiplicative,
 linearly bounded, and left continuous; its Fubini--Study function is a
 continuous function on $(\mathbf P^1)^{\An}$, indeed a
 psh function for $c_1(\cO(1))$ by the curve criterion
 \cite[Example~4.13]{BJ22}: the two nonconstant rays have initial slopes
 $-\tfrac12$ and $-\tfrac12$, all other initial slopes vanish, and
 $1+(-\tfrac12)+(-\tfrac12)=0$, while every ray restriction is convex.
 Its Duistermaat--Heckman measure exists.  This is the precise sense in
 which filtrations are a strictly larger category than test configurations.
\end{enumerate}

\emph{Its Duistermaat--Heckman measure.}  The concave transform on
$\Delta=[0,1]$ is $G_{\mathcal F}(u)=u(1-u)$, so by
\eqref{eq:ch33-dh-is-pushforward} the measure is the push-forward of
Lebesgue measure on $[0,1]$ by $u\mapsto u(1-u)$.  For
$0<\tau<1/4$ the superlevel set $\{G_{\mathcal F}>\tau\}$ is the
interval of length $\sqrt{1-4\tau}$ centred at $1/2$; hence, by
\Cref{prop:dh-tail},
\begin{equation}
\label{eq:ch33-harmonic-dh}
 \DHm(\mathcal F)
 =\frac{2}{\sqrt{1-4\tau}}\,\mathbf 1_{[0,1/4)}(\tau)\,\rd\tau ,
\end{equation}
of total mass $\int_0^{1/4}2(1-4\tau)^{-1/2}\rd\tau=1$, with no atom and
with unbounded density at the right endpoint.  Its first moment is
$\int_0^1u(1-u)\rd u=1/6$.  Note that the density is \emph{increasing};
this is consistent with \Cref{cor:ch33-dh-structure}\,(2), which for
$n=1$ asserts concavity of $\tau\mapsto\vol\{G_{\mathcal F}\geq\tau\}$,
and $\tau\mapsto\sqrt{1-4\tau}$ is indeed concave.
\end{example}

\begin{remark}
\label{rem:ch33-harmonic-reverse}
The measure \eqref{eq:ch33-harmonic-dh} has a continuous density that is not
piecewise constant.  By \cite[Theorem~5.10]{BHJ17}, a finitely generated
$\mathbf Z$-filtration on a normal projective curve has piecewise constant
density for the absolutely continuous part of its limit measure.
Restricting the harmonic filtration to integer levels replaces its monomial
weights by their floors.  The normalized jumping lists then differ by
less than $1/m$, so the two filtrations have the same DH measure.
Consequently the integer-level Rees algebra in
\Cref{ex:ch33-harmonic-not-fg} cannot be finitely generated either.
This gives a second proof of that algebraic assertion without extending
the cited theorem to arbitrary real-weight filtrations.  Although far
from a complete invariant, the DH measure detects this obstruction.
\end{remark}

\subsection{Toric test configurations}
\label{sec:ch33-example-toric}

On a toric variety every object in the dictionary is a convex-geometric
object, and the whole chapter can be read off from a polytope and a concave
function on it.  We use the toric conventions of
\Cref{sec:toric-convex} and \eqref{eq:appJ-toric-coordinate}: the
tropical coordinate of a point of the dense torus is
$u$ with $\langle m,u\rangle=-\log|\chi^m|$, and the monomial
semivaluation $v_u$ satisfies $v_u(\chi^m)=\langle m,u\rangle$.

\begin{theorem}
\label{thm:ch33-toric}
Let $X_\Sigma$ be a proper toric variety over $k$ of dimension $n$,
with lattice $N$, dual lattice $M$, and let $L$ be an ample toric
line bundle with polytope $P\subset M_{\mathbf R}$, so that
$R_m=\bigoplus_{u\in mP\cap M}k\,\chi^u$ and
$V=(L^n)=n!\vol(P)$.  Let $g:P\to\mathbf R$ be continuous and concave.
Define
\begin{equation}
\label{eq:ch33-toric-filtration}
 \chi_m\Bigl(\sum_{u\in mP\cap M}c_u\chi^u\Bigr)
 :=\min\{m\,g(u/m)\mid c_u\ne0\} .
\end{equation}
Then:
\begin{enumerate}
\item $\chi$ is a multiplicative, linearly bounded graded norm, and the
 corresponding filtration is spanned by characters in each degree;
\item its Duistermaat--Heckman measure is
\begin{equation}
\label{eq:ch33-toric-dh}
 \DHm(\mathcal F)=g_\ast\Bigl(\frac{\lambda_P}{\vol(P)}\Bigr),
 \qquad
 \int\lambda\,\rd\DHm(\mathcal F)
 =\frac1{\vol(P)}\int_Pg(x)\,\rd x ;
\end{equation}
\item for every $u\in N_{\mathbf R}$,
\begin{equation}
\label{eq:ch33-toric-fs}
 \phi_\chi(v_u)
 =\sup_{x\in P}\bigl\{g(x)-\langle x,u\rangle\bigr\}
   -\check h_P(u),
 \qquad
 \check h_P(u)=\sup_{x\in P}\langle x,-u\rangle ;
\end{equation}
\item $\mathcal F$ is of finite type if and only if $g$ is affine on
 the cells of a finite rational polyhedral subdivision of $P$.
 Its integer-rounded filtration, with weights
 $\widetilde\chi_m=\lfloor\chi_m\rfloor$, is finitely generated if and
 only if, in addition, the affine data are rational.  The rounded
 filtration has the same function $\phi_\chi$ and DH measure; it is
 the integral-weight presentation used to produce an algebraic test
 configuration.  The original real weights need not be integral on any
 Veronese section ring;
\item if $g\geq0$ is piecewise affine with rational data, then the
 unnormalized first moment is
\begin{equation}
\label{eq:ch33-toric-intersection}
 \int\lambda\,\rd\DHm^{\mathrm u}(\mathcal F)
 =n!\int_Pg(x)\,\rd x
 =n!\,\vol(Q),
 \qquad
 (\overline\cL^{\,n+1})=(n+1)!\,\vol(Q),
\end{equation}
 where $Q=\{(x,\tau)\in P\times\mathbf R\mid 0\leq\tau\leq g(x)\}$.
\end{enumerate}
\end{theorem}

\begin{proof}
(1)  Concavity of $g$ means that
$\Psi(u,m):=m\,g(u/m)$ is concave and positively homogeneous of degree one
on the cone $C:=\operatorname{Cone}(P\times\{1\})\subset M_{\mathbf
R}\times\mathbf R$, hence superadditive; the argument of
\Cref{ex:ch33-harmonic-not-fg} then gives
\eqref{eq:graded-norm-inequalities}, using that the coefficient of a
character in a product is a sum of products of coefficients.  Boundedness of
$g$ on the compact $P$ gives linear boundedness.  The filtered pieces
are spanned by characters by construction.

(2)  Let $f$ be continuous and bounded on $\mathbf R$.  Then
\[
 \int f\,\rd\nu_m
 =\frac1{\#(mP\cap M)}\sum_{u\in mP\cap M}f\bigl(g(u/m)\bigr),
\]
a Riemann sum for $\vol(P)^{-1}\int_Pf(g(x))\rd x$ over the lattice
$m^{-1}M$.  Precisely, $m^{-n}\#(mP\cap M)\to\vol(P)$ and
$m^{-n}\sum_{u\in mP\cap M}h(u/m)\to\int_Ph$ for every continuous $h$
on $P$, by uniform continuity of $h$ and the fact that the cells
$u/m+m^{-1}[0,1)^n$ tile a neighbourhood of $P$ with volume
$m^{-n}$; apply this to $h=f\circ g$.  Hence
$\nu_m\to g_\ast(\lambda_P/\vol(P))$ weakly, which is
\eqref{eq:ch33-toric-dh}; the moment formula is \Cref{lem:dh-moments}.

(3)  Let $u\in N_{\mathbf R}$ and let $\sigma$ be a cone of the normal
fan containing $u$, with corresponding vertex $u_\sigma$ of $P$, so
that $\langle u_\sigma,u\rangle=\min_{x\in P}\langle x,u\rangle$.  On the
affine chart $U_\sigma$, which contains the center of $v_u$, the
line bundle $mL$ is trivialized by $\chi^{mu_\sigma}$, and the section
$\chi^{u'}$ becomes the function $\chi^{u'-mu_\sigma}$; hence
$v_u(\chi^{u'})=\langle u'-mu_\sigma,u\rangle$.  Therefore
\[
 \frac{\chi_m(\chi^{u'})-v_u(\chi^{u'})}{m}
 =g(u'/m)-\bigl\langle u'/m-u_\sigma,u\bigr\rangle ,
\]
and taking the supremum over $m$ and over $u'\in mP\cap M$, which is a
supremum over a dense subset of $P$ of a continuous function, gives
\[
 \phi_\chi(v_u)=\sup_{x\in P}\{g(x)-\langle x,u\rangle\}
   +\langle u_\sigma,u\rangle .
\]
Since $\langle u_\sigma,u\rangle=\min_{x\in P}\langle x,u\rangle
=-\check h_P(u)$, this is \eqref{eq:ch33-toric-fs}.  As in the previous
examples, the supremum in \eqref{eq:norm-fs-function} may be restricted to
characters, because both $\chi_m$ and $v_u$ are minima over
the character support.

(4)  Suppose first that $\mathcal F$ is of finite type.
Replace its finite generating sections by an adapted character basis
in the same degree $r$.  This preserves generation: each old generator
is a linear combination of characters whose weights are at least its
weight, and expansion of products gives one inclusion; multiplicativity
gives the reverse inclusion.  The finite-maximum formula then yields
\[
 g^\ast(u):=\sup_{x\in P}\{g(x)-\langle x,u\rangle\}
 =\max_{a\in rP\cap M}
       \{\chi_r(\chi^a)/r-\langle a/r,u\rangle\}.
\]
The concave Fenchel--Moreau formula recovers $g$ as the upper roof
of the convex hull of the finitely many points
$(a/r,\chi_r(\chi^a)/r)$.  The projections of its upper faces are
convex hulls of rational points $a/r$.  They give a finite rational
polyhedral subdivision on which $g$ is affine.  This argument uses
characters; the valuation of an arbitrary section is a minimum of
affine terms and need not itself be affine on a fan cone.

Conversely, suppose $g$ is affine on such a subdivision.  Each cone
over a cell has a finite Hilbert basis of its lattice semigroup.
Assign each Hilbert-basis section its real weight $\Psi(u,m)$.
Every lattice point decomposes in one cell, where $\Psi$ is linear,
so its exact weight is the sum of the generator weights.  These
finitely many weighted sections, possibly of different degrees,
generate the filtration.  Pass to a common generating Veronese of
the polynomial algebra on these generators: partitioning a monomial
into the new degree-one monomials preserves its total real weight.
Their images satisfy \Cref{def:ch33-finite-type}.  Thus the original
real filtration is of finite type.

The integer-level Rees algebra has semigroup
\[
 S=\{(u,m,\lambda)\in M\times\mathbf Z\times\mathbf Z:
                 u\in mP,\ \lambda\leq\Psi(u,m)\}.
\]
Its finite generation is equivalent to rational polyhedrality of the
hypograph by \Cref{prop:ch33-monomial-fg-criterion}.  This is exactly
piecewise affinity with rational affine data.

Finally, rounding down preserves multiplicativity since
$\lfloor a\rfloor+\lfloor b\rfloor\leq\lfloor a+b\rfloor$.
It preserves $\phi_\chi$: one inequality follows from
$\widetilde\chi\leq\chi$, and the other follows by applying
$\widetilde\chi_{jm}(s^j)\geq j\chi_m(s)-1$ and sending $j\to\infty$.
The ordered jumping lists differ by at most one before normalization,
hence by at most $1/m$ afterwards.  Their uniformly supported empirical
measures have the same weak limit.  This justifies using the rounded
filtration for the algebraic construction without identifying its
individual weights with the original real ones.

(5)  The first equality is \eqref{eq:ch33-toric-dh} multiplied by
$V=n!\vol(P)$; the second is Fubini for the region under the graph of
$g$.  The last is
\eqref{eq:ch33-first-moment-intersection} for the associated
integer-rounded toric configuration.  Thus Fubini's volume of the
region under $g$ and the algebraic first-moment identity give the
displayed intersection formula.
\end{proof}

\begin{example}
\label{ex:ch33-toric-p2}
Let $X=\mathbf P^2$, $L=\cO(1)$, so that
$P=\Delta=\operatorname{conv}\{(0,0),(1,0),(0,1)\}$,
$\vol(\Delta)=1/2$, $n=2$ and $V=2!\cdot\tfrac12=1=(\cO(1)^2)$.  Fix
$c\in\mathbf Q$ with $0<c<1/2$ and put
\[
 g(x_1,x_2)=\min\{x_1,\ x_2,\ c\} ,
\]
which is concave, continuous, and nonnegative, with values in $[0,c]$.
Identify $R_m=H^0(\mathbf P^2,\cO(m))$ with the span of the monomials
$X^aY^bZ^{m-a-b}$, $(a,b)\in m\Delta\cap\mathbf Z^2$.

\emph{The filtration and its base ideals.}  By
\eqref{eq:ch33-toric-filtration},
\[
 \chi_m(X^aY^bZ^{m-a-b})=\min\{a,\ b,\ mc\} ,
\]
so that for $0\leq\lambda\leq mc$ the space
$\mathcal F^\lambda R_m$ is spanned by the monomials with
$a\geq\lambda$ and $b\geq\lambda$, and
$\mathcal F^\lambda R_m=0$ for $\lambda>mc$.  For
$c=1/4$ and $m=4$ this reads
\[
\begin{gathered}
 \mathcal F^0R_4=R_4,\qquad
 \mathcal F^2R_4=0 ,\\
 \mathcal F^1R_4=\langle X^aY^bZ^{4-a-b}\mid a,b\geq1\rangle
 =H^0\bigl(\mathbf P^2,\cO(4)\bigr)\cap(XY),
\end{gathered}
\]
with base ideals $\mathfrak b_{4,0}=\cO_X$,
$\mathfrak b_{4,1}=\cO_X(-D_X-D_Y)$, where $D_X=\{X=0\}$ and
$D_Y=\{Y=0\}$.  The flag ideal of
\eqref{eq:ch33-degree-m-flag} is therefore
\begin{equation}
\label{eq:ch33-toric-flag}
 \mathfrak a_{(4)}
 =\cO_X+\cO_X(-D_X-D_Y)\,\varpi^{-1},
 \qquad
 \tfrac14\phi_{\mathfrak a_{(4)}}
 =\tfrac14\max\bigl\{0,\ \log|\cO(-D_X-D_Y)|+1\bigr\} .
\end{equation}

\emph{The Fubini--Study function, twice.}  By \eqref{eq:ch33-toric-flag},
$\phi_\chi(v)=\tfrac14\max\{0,1-v(D_X)-v(D_Y)\}$.  On the other hand
\eqref{eq:ch33-toric-fs} gives, for $u$ in the cone
$\sigma_0=\mathbf R_{\geq0}^2$ of the normal fan, where
$\check h_\Delta(u)=0$ and $v_u(D_X)=u_1$, $v_u(D_Y)=u_2$,
\[
 \phi_\chi(v_u)
 =\sup_{x\in\Delta}\bigl\{\min\{x_1,x_2,\tfrac14\}
      -x_1u_1-x_2u_2\bigr\}
 =\tfrac14\max\{0,\ 1-u_1-u_2\} .
\]
For instance $\phi_\chi(v_{(0,0)})=1/4$,
$\phi_\chi(v_{(1,0)})=0$, and
$\phi_\chi(v_{(1/4,1/4)})=1/8$, the last being attained at
$x=(1/4,1/4)$.  The two computations agree, which checks
\eqref{eq:ch33-flag-vs-norm} in a two-dimensional example.  A point of the
cone $\operatorname{Cone}((-1,-1),(0,1))$, say $u=(-1,0)$, has
$\check h_\Delta(u)=1$ and $\sup_{x}\{g(x)+x_1\}=1$, attained at
$x=(1,0)$; so $\phi_\chi(v_{(-1,0)})=0$, in agreement with the
flag-ideal formula, since in the corresponding chart $D_X$ is not
visible while $D_Y$ has local equation $\chi^{(-1,1)}$, whence
$v_u(D_X)=0$ and $v_u(D_Y)=1$.

\emph{The Duistermaat--Heckman measure.}  The pair $(\Delta,g)$ is exactly
the one computed in \Cref{ex:dh-simplex}, where
$\DHm_{\Delta,g}
=2(1-2\tau)\mathbf 1_{[0,c)}\rd\tau+\tfrac12(1-2c)^2\delta_c$.  Dividing by
$\vol(\Delta)=1/2$ gives the probability normalization used here:
\begin{equation}
\label{eq:ch33-toric-p2-dh}
 \DHm(\mathcal F)
 =4(1-2\tau)\,\mathbf 1_{[0,c)}(\tau)\,\rd\tau
   +(1-2c)^2\,\delta_c ,
\end{equation}
of total mass $4(c-c^2)+1-4c+4c^2=1$.  The atom at $\tau=c$ is
$\vol\{g=c\}/\vol(\Delta)$, which is the truncation locus of $g$; the
density $4(1-2\tau)$ is decreasing.  This is consistent with
\Cref{cor:ch33-dh-structure}\,(2) but is not forced by it: for $n\geq2$
concavity of $\tau\mapsto\vol\{G\geq\tau\}^{1/n}$ leaves the monotonicity
of the density open, as $W(\tau)=(1-\tau)^{1/3}$ with $n=2$ shows ---
$W$ is concave, yet $W^2$ has density $\tfrac23(1-\tau)^{-1/3}$, which
increases.  What forces it here is the stronger fact that
$\vol\{g\geq\tau\}^{1/2}=2^{-1/2}(1-2\tau)$ is \emph{affine}, so that
$\vol\{g\geq\tau\}$ is a quadratic with nondecreasing derivative, whose negative is the
nonincreasing density.

\emph{Moments and intersection numbers.}  By the layer-cake formula
$\int_\Delta g=\int_0^c\vol\{g>\tau\}\rd\tau
=\int_0^c\tfrac12(1-2\tau)^2\rd\tau=\tfrac1{12}\bigl(1-(1-2c)^3\bigr)$, so
\[
 \int\lambda\,\rd\DHm(\mathcal F)=\frac16\bigl(1-(1-2c)^3\bigr),
 \qquad
 (\overline\cL^{\,3})
 =3\int\lambda\,\rd\DHm^{\mathrm u}(\mathcal F)
 =\frac12\bigl(1-(1-2c)^3\bigr),
\]
using $V=1$ and \eqref{eq:ch33-toric-intersection}.  For $c=1/4$ this is
$\int\lambda\,\rd\DHm=\tfrac7{48}$ and
$(\overline\cL^{\,3})=\tfrac7{16}$.  As $c\to0$ both tend to $0$, and
the configuration degenerates to the trivial one; as $c\to1/2$ the atom
disappears and $\DHm$ becomes $4(1-2\tau)\mathbf 1_{[0,1/2)}\rd\tau$.

\emph{Finite generation.}  The function $g=\min\{x_1,x_2,c\}$ is a minimum
of three rational affine functions, so
\Cref{thm:ch33-toric}\,(4) applies: $\mathcal F$ is rational and finitely
generated, and by \Cref{thm:ch33-fg-filtration-tc} it comes from an
integrally closed semiample test configuration.  For $c=1/4$, a representative is obtained by blowing up
\eqref{eq:ch33-toric-flag}.  If $\mathfrak a_{(4)}\cO_{\cX}=\cO_{\cX}(D)$,
its polarization is
$\cL=\mu^\ast(4L)_{\triv}+D$ divided by $4$.
The tautological vertical divisor alone is not a polarization extending
$4L$.  For other rational $c$, use the corresponding sufficiently
divisible degree-specific flag ideal.  Replacing $g$ by the
continuous nonpolyhedral concave function
$\min\{c,\sqrt{x_1x_2}\}$ destroys finite generation for the reason
isolated in \Cref{ex:ch33-harmonic-not-fg}; the general identities
\eqref{eq:ch33-toric-dh} and \eqref{eq:ch33-toric-fs} still apply, with the
new $g$.
\end{example}

\section{The deferred geodesic dictionary}
\label{sec:ch33-geodesic-deferred}

Over $\mathbf C$, filtrations also enter the construction of weak
geodesic rays, but that analytic correspondence is not part of the present
chapter.  Its three steps are proved later: filtrations give test curves in
\Cref{thm:rf-filtration-test-curve}, test configurations give the
Phong--Sturm rays of \cite{PS07} in \Cref{thm:rf-phong-sturm}, and the
available energy-slope statement is isolated in
\Cref{thm:rf-energy-slope}.  Those results require partial Bergman kernels,
maximal envelopes, and the Ross--Witt Nystr\"om correspondence.

The algebraic data needed there have already been constructed here: the
filtered section spaces, the function $\phi_\chi$, and the empirical
measures $\nu_m$.  No result in
\Crefrange{chap:bj-fs-pl-psh}{chap:bj-stability} uses the complex
geodesic construction.

\begin{warning}{Nothing in Part IV depends on the complex side}
\label{warn:ch33-no-complex-input}
Every statement proved in this chapter is algebraic and valid over an
arbitrary algebraically closed trivially valued field, of arbitrary
characteristic.  In particular
\Cref{thm:dh-limit-filtration,thm:ch33-dh-weight-count} do not use complex
geometry, and the energy functionals of \Cref{chap:bj-energy} are
constructed from the non-Archimedean side alone.  A reader interested only
in Part IV may skip the remainder of this section.
\end{warning}

\subsection{The erratum boundary}
\label{sec:ch33-erratum-boundary}

Asymptotic slope statements---identifications of the limit as
$t\to\infty$ of $t^{-1}$ times a complex energy functional along a ray
with a non-Archimedean functional---must be handled with care, because part
of the published literature on them was retracted.  The precise dependency
ledger is \Cref{thm:rf-erratum-audit}, which we adopt verbatim as a standing
constraint.

\begin{warning}{No erratum-affected input}
\label{warn:ch33-erratum-firewall}
No statement of this chapter uses, directly or transitively, any of the
results listed in \eqref{eq:rf-quarantined-results}, nor any reformulation
of one of them under another name; the prohibition of
\Cref{warn:rf-no-surrogate} is in force here as well.  The energy-related
asymptotic identities used above are the purely algebraic moment formulas: the first-moment formula
\eqref{eq:ch33-first-moment-intersection} comes from
\cite[Proposition~3.12(iii)]{BHJ17}, and higher moments come from
Corollary~3.4 there.  This is a different paper, concerned with
Duistermaat--Heckman measures of test configurations over an arbitrary
field, and unaffected by the erratum.  No statement of this chapter asserts
that a complex energy slope equals a non-Archimedean energy; such
statements appear only in \Cref{chap:rays-filtrations}, each with its own
post-erratum source.
\end{warning}

\section{Dictionary and provenance}
\label{sec:ch33-dictionary}

The table records the principal maps, identities, and equivalences proved
above.  Only rows marked $\longleftrightarrow$ are equivalences; in
particular the Duistermaat--Heckman measure is a coarse invariant rather
than a reconstruction theorem.

\begin{center}
\small
\renewcommand{\arraystretch}{1.18}
\begin{tabularx}{\textwidth}{@{}p{.25\textwidth}>{\raggedright\arraybackslash}X p{.24\textwidth}@{}}
\toprule
Data & Relation or invariant & Result\\
\midrule
Test configuration $\cX$
 &$\longrightarrow$ integrally closed representative
 &\Cref{prop:ch33-integral-closure-exists}\\
Compactification $\overline\cX\to\mathbf P^1$
 &$\longrightarrow$ intersection invariants
 &\Cref{thm:ch33-intersection-invariance}\\
Flag ideal $\mathfrak a$
 &$\longrightarrow$ relatively integrally closed blowup
 &\Cref{thm:ch33-flag-tc-equivalence}\\
Globally generated vertical divisor $D$
 &$\longrightarrow$ function $\phi_D=\phi_{\mathfrak a}$
 &\Cref{thm:ch33-flag-tc-equivalence}\\
Semiample $(\cX,\cL)$
 &$\longleftrightarrow$ saturated f.g.\ integral-weight filtration
 &\Cref{thm:ch33-fg-filtration-tc}\\
Ample central-fibre representation, $m\gg0$
 &$\longleftrightarrow\operatorname{gr}_{\mathcal F}R_m$
 &\Cref{prop:ch33-central-fibre-weights}\\
Raw base change of degree $d$
 &$\longrightarrow\chi_m\mapsto d\chi_m$
 &\Cref{prop:ch33-transformations}\\
Character twist
 &$\longrightarrow\chi_m\mapsto\chi_m+cm$
 &\Cref{prop:ch33-transformations}\\
Filtered subspaces
 &$\longrightarrow\DHm(\mathcal F)=\dfrac{n!}{V}
   (G_{\mathcal F})_\ast\lambda_\Delta$
 &\Cref{thm:ch33-dh-concave-transform}\\
First moment
 &$V\!\displaystyle\int\lambda\,\rd\DHm
   =\dfrac{(\overline\cL^{\,n+1})}{n+1}$
 &\Cref{thm:ch33-dh-weight-count}\\
\bottomrule
\end{tabularx}
\end{center}

Three boundaries of the dictionary should be kept in mind.  Filtrations that
are not of finite type have no semiample model realizing those filtered pieces
(\Cref{ex:ch33-harmonic-not-fg}); the section-supremum function determines the equivalence class of a
semiample test configuration with fixed generic polarization and
linearization (\Cref{prop:tc-valuative-invariance}); and the
Duistermaat--Heckman measure, while a genuine invariant, is far from a
complete one, since it depends only on the jumping numbers and not on the
position of the filtered subspaces.

\section*{Source note}

The test-configuration and Fubini--Study dictionary is
\cite[\S\S1--2]{BJ22}; the filtration and Duistermaat--Heckman viewpoint is
developed in \cite{BHJ17}.  The definitions above retain the integral
closure and the linearization, since forgetting either makes the dictionary
noninjective.

In more detail: \Cref{def:bj-test-configuration} and the structure results
of \Cref{sec:ch33-tc-structure} follow \cite[\S1.2]{BJ22}, which in turn
rests on \cite[\S2]{BHJ17}; reducedness of the total space is
\cite[Proposition~2.6]{BHJ17}.  Integrally closed test configurations, the
vertical Serre criterion \Cref{thm:ch33-integrally-closed-criterion}, and
the polar-scheme argument in \Cref{lem:ch33-effectivity} are
\cite[\S1.4]{BJ22}; the notion goes back to the ``partially normal''
configurations of \cite{Od13}, and its use in K-stability to
\cite{Od13,Od13b}.  The correspondence between components of the central
fibre and divisorial valuations is \cite[Theorem~4.6]{BHJ17}, quoted as
\cite[Lemma~1.20]{BJ22}.  Flag ideals are due to \cite{Od13}; their
semiring calculus \eqref{eq:flag-semiring} and the associated PL functions
are \cite[\S2.1]{BJ22}, and Rees valuations are
\cite[Definition~2.10]{BJ22}.  The flag-ideal construction in
\Cref{thm:ch33-flag-tc-equivalence} uses
\cite[Lemmas~2.8--2.9 and Proposition~2.25]{BJ22};
\cite[Lemma~2.12]{BJ22} supplies the Rees-valuation realization, while the
unique ample integrally closed representative is
\cite[Proposition~2.30]{BJ22}.  The full test-configuration/PL
correspondence \cite[Theorem~2.31]{BJ22} is deferred to
\Cref{chap:bj-fs-pl-psh}.  Finally,
\Cref{lem:ch33-projective-blowup-presentation} uses
\cite[Lemma~1.8]{BJ22} together with relative $\Proj$ reconstruction.
The uniqueness theorem does not identify an arbitrary finitely generated
filtration with its saturated integral closure.  That quotient is stated explicitly in
\Cref{thm:ch33-fg-filtration-tc}.

Filtrations of section rings entered K-stability through
\cite{Sze15}, whose appendix by Boucksom relates general filtrations to
test configurations; the systematic treatment in the form used here is
\cite[\S2]{BHJ17}.  More precisely, the ample finite-generation
correspondence is \cite[Proposition~2.15]{BHJ17}, the ample model of a
semiample configuration is \cite[Proposition~2.17]{BHJ17}, and the
degree-specific flag ideals used in
\Cref{prop:ch33-base-ideal-flag} are
\cite[Definition~2.19 and Proposition~2.21]{BHJ17}.  The
Duistermaat--Heckman measure of a test
configuration, its moments, and the intersection formula
\eqref{eq:ch33-first-moment-intersection} are \cite[\S3]{BHJ17},
specifically Proposition~3.12(iii) for the original compactification's
first moment and Corollary~3.4 for the auxiliary fibrations' higher moments.  The existence of the limit for an
arbitrary linearly bounded multiplicative filtration, and its description as
the push-forward of Lebesgue measure on an Okounkov body by a concave
transform, is \cite[Theorem~A]{BC11}; upper semicontinuity of the concave
transform is built into \cite[Definition~1.8]{BC11}.  This rests on the Okounkov-body
construction of \cite{Oko03,LM09,KK12}; the relation between test
configurations and Okounkov bodies was first exploited in \cite{WN12}.  The
convexity properties recorded in \Cref{cor:ch33-dh-structure} are the
general convex-geometric statements of \Cref{sec:dh-preview}, in particular
Brunn--Minkowski concavity of the tail volumes.  Toric test configurations in the form of
\Cref{thm:ch33-toric} go back to \cite{Don02}.  Metric properties of the space
of graded norms, including the distances $d_p$ for which
\Cref{prop:ch33-dh-continuity} is one endpoint, are developed in
\cite{BJKS1} and used in \Cref{chap:bj-stability}.  Finite generation for
the filtrations that compute stability thresholds is the theorem of
\cite{LXZ22}; nothing in this chapter depends on it.

Two normalizations deserve a final word.  The Duistermaat--Heckman measure
of \Cref{thm:dh-limit-filtration} is a probability measure, matching
\cite{BHJ17} and \Cref{thm:rf-dh-quantization}, while
\Cref{sec:appJ-rays} and \Cref{chap:transcendental-okounkov-dh} use the
unnormalized measure of mass $V$; the conversion is
\eqref{eq:ch33-dh-unnormalized}.  And the sign of the filtration weight is
fixed by \Cref{def:ch33-tc-filtration} and
\Cref{rem:ch33-weight-sign}, not by an appeal to a one-parameter subgroup;
a source that lets the subgroup act on points rather than on homogeneous
coordinates obtains the reflected measure.
